\documentclass[11pt]{article}
\usepackage[margin=1in]{geometry}
\usepackage[T1]{fontenc}
\usepackage[utf8]{inputenc}
\usepackage[french]{babel}
\usepackage{amsmath,amssymb,amsthm}
\usepackage[colorlinks=true,linkcolor=blue,urlcolor=blue]{hyperref}
\newtheorem{problem}{Problème}
\newenvironment{refs}{\par\small\noindent\emph{Références :}\begin{itemize}\setlength\itemsep{0pt}}{\end{itemize}\normalsize}
\title{Hors de la Caverne \\ \Large Problèmes ouverts et frontières de la recherche}
\author{}
\date{}
\begin{document}
\maketitle
\noindent\emph{Des problèmes, pas de réponses.} Chaque problème ci-dessous est posé
sans solution. Les références indiquent où trouver les connaissances nécessaires.
\medskip


\section*{Les problèmes du prix du millénaire}

\begin{problem}[{L'hypothèse de Riemann --- Recherche}]
\textbf{OPEN-01.} \textbf{Problème.} La fonction zêta de Riemann est définie pour
$ \operatorname{Re}(s) > 1 $ par
\[ \zeta(s) = \sum_{n=1}^{\infty} \frac{1}{n^s} \]
puis prolongée au plan complexe tout entier (privé d'un pôle en $ s = 1 $) par
prolongement analytique. Elle s'annule aux entiers pairs négatifs $ -2, -4, -6, \dots $ ;
ce sont les zéros triviaux. \textbf{Démontrez que tout autre zéro de $ \zeta $ vérifie
$ \operatorname{Re}(s) = \tfrac{1}{2} $.} Statut : ouvert depuis 1859.
\end{problem}

\noindent Riemann a introduit cette question dans un mémoire de huit pages daté
de 1859, le seul article qu'il ait jamais écrit sur la théorie des nombres, en expliquant
comment compter les nombres premiers. Sa formule explicite exprime l'erreur dans la
fonction de comptage des nombres premiers comme une somme sur les zéros de $ \zeta $ ---
si bien que la position des zéros contrôle littéralement le degré d'irrégularité de la
répartition des nombres premiers. L'hypothèse affirme que les nombres premiers sont aussi
réguliers qu'ils pourraient l'être. Hilbert en a fait le huitième de ses problèmes de 1900
et aurait déclaré que, s'il se réveillait après cinq cents ans, sa première question serait
de savoir si elle avait été démontrée. Plus de $ 10^{13} $ zéros ont été vérifiés et tous
se trouvent sur la droite critique ; c'est un indice, pas une démonstration, et les
mathématiques se sont déjà laissé tromper par des indices numériques plus massifs encore
(voir le nombre de Skewes). Des centaines de théorèmes publiés commencent par \og{} supposons
l'hypothèse de Riemann \fg{}.

\begin{refs}
\item Énoncé officiel : Clay Mathematics Institute --- Riemann Hypothesis (\url{https://www.claymath.org/millennium/riemann-hypothesis/})
\item Contexte : Hypothèse de Riemann --- Wikipédia (\url{https://fr.wikipedia.org/wiki/Hypoth%C3%A8se_de_Riemann})
\item Lecture : John Derbyshire, \textit{Prime Obsession} (2003) --- l'histoire, pour un lecteur non spécialiste
\item Lecture : H. M. Edwards, \textit{Riemann's Zeta Function} --- contient une traduction du mémoire de 1859
\item Lecture : Marcus du Sautoy, \textit{The Music of the Primes} (2003)
\end{refs}

\bigskip

\begin{problem}[{P contre NP --- Recherche}]
\textbf{OPEN-02.} \textbf{Problème.} $ \mathrm{P} $ est la classe des problèmes de
décision que résout une machine de Turing déterministe en un temps polynomial en la
longueur de l'entrée. $ \mathrm{NP} $ est la classe dont les instances \og{} oui \fg{}
possèdent des certificats vérifiables en temps polynomial. \textbf{Décidez si
$ \mathrm{P} = \mathrm{NP} $.} Statut : ouvert. On croit très largement que
$ \mathrm{P} \neq \mathrm{NP} $, et aucune démonstration n'est connue dans l'une ou
l'autre direction.
\end{problem}

\noindent La question a été posée sous une forme essentiellement moderne par
Stephen Cook en 1971 et, indépendamment, par Leonid Levin ; une lettre de 1956 de Kurt
Gödel à John von Neumann l'avait déjà formulée dans un autre langage. C'est la question
de savoir si trouver une solution est fondamentalement plus difficile que la reconnaître
--- si la créativité peut être automatisée. Des milliers de problèmes naturels
(l'ordonnancement, les modèles de repliement des protéines, la conception de circuits, le
sudoku sur des grilles $ n \times n $) sont NP-complets, ce qui signifie qu'un
algorithme polynomial pour l'un quelconque d'entre eux ferait s'effondrer la classe tout
entière. Deux barrières expliquent pourquoi les progrès sont si difficiles : la
relativisation (Baker--Gill--Solovay, 1975) montre que toute démonstration traitant le
calcul comme une boîte noire est vouée à l'échec, et les démonstrations naturelles
(Razborov--Rudich, 1994) montrent qu'une vaste classe de techniques combinatoires ne peut
pas fonctionner, sauf à rendre la cryptographie impossible.

\begin{refs}
\item Énoncé officiel : Clay Mathematics Institute --- P vs NP (\url{https://www.claymath.org/millennium/p-vs-np/})
\item Contexte : Problème P ≟ NP --- Wikipédia (\url{https://fr.wikipedia.org/wiki/Probl%C3%A8me_P_%E2%89%9F_NP})
\item Lecture : Lance Fortnow, \textit{The Golden Ticket: P, NP, and the Search for the Impossible} (2013)
\item Lecture : Sanjeev Arora \& Boaz Barak, \textit{Computational Complexity: A Modern Approach}
\item Voir aussi : Mathématiques discrètes et informatique (\url{applied-discrete-cs.html})
\end{refs}

\bigskip

\begin{problem}[{Existence et régularité pour Navier--Stokes --- Recherche}]
\textbf{OPEN-03.} \textbf{Problème.} Pour un écoulement incompressible dans
$ \mathbb{R}^3 $, de champ de vitesse $ u $ et de pression $ p $, les équations de
Navier--Stokes s'écrivent
\[ \partial_t u + (u \cdot \nabla) u = -\nabla p + \nu \Delta u, \qquad \nabla \cdot u = 0 . \]
\textbf{Démontrez que, pour tout champ de vitesse initial lisse et à décroissance rapide, il
existe une solution lisse définie pour tout temps $ t > 0 $} --- ou exhibez une donnée
initiale dont la solution explose en temps fini. Statut : ouvert en dimension trois.
\end{problem}

\noindent Les équations ont été écrites par Claude-Louis Navier en 1822 et
solidement établies par George Gabriel Stokes en 1845 ; on s'en sert tous les jours pour
concevoir des avions et prévoir le temps. Personne ne sait si elles admettent toujours des
solutions. Jean Leray a démontré en 1934 que des solutions faibles existent toujours, mais
il n'a pu montrer ni leur unicité ni leur régularité ; cet écart n'a jamais été comblé.
L'enjeu physique est la description mathématique de la turbulence. Notez l'asymétrie, dite
honnêtement : en dimension deux, la régularité globale est un \textit{théorème}, démontré par
Ladyjenskaïa ; c'est précisément la troisième dimension, où les lignes de tourbillon
peuvent s'étirer, qui résiste.

\begin{refs}
\item Énoncé officiel : Clay Mathematics Institute --- Navier--Stokes Equation (\url{https://www.claymath.org/millennium/navier-stokes-equation/})
\item Contexte : Solutions des équations de Navier-Stokes --- Wikipédia (\url{https://fr.wikipedia.org/wiki/Solutions_des_%C3%A9quations_de_Navier-Stokes})
\item Lecture : Terence Tao, \og{} Finite time blowup for an averaged three-dimensional Navier-Stokes equation \fg{} (2016), arXiv:1402.0290 (\url{https://arxiv.org/abs/1402.0290}) --- pourquoi le problème est véritablement délicat
\item Voir aussi : Équations aux dérivées partielles (\url{applied-pde.html})
\end{refs}

\bigskip

\begin{problem}[{Existence de Yang--Mills et déficit de masse --- Recherche}]
\textbf{OPEN-04.} \textbf{Problème.} \textbf{Démontrez que, pour tout groupe de jauge simple
compact $ G $, une théorie quantique de Yang--Mills non triviale existe sur
$ \mathbb{R}^4 $ et satisfait les axiomes de Wightman, et qu'elle possède un déficit de
masse $ \Delta > 0 $} --- c'est-à-dire que tout état autre que le vide a une énergie au
moins égale à $ \Delta $. Statut : ouvert. Cela suppose d'abord de construire la théorie
de façon rigoureuse, ce qui est déjà non résolu.
\end{problem}

\noindent La théorie de Yang--Mills, écrite par Chen Ning Yang et Robert Mills en
1954, est l'ossature mathématique du modèle standard de la physique des particules. Les
physiciens calculent avec elle tous les jours et ses prédictions s'accordent avec
l'expérience avec une précision extraordinaire ; l'embarras, c'est que personne n'a
construit la théorie comme objet mathématique en dimension quatre. Le déficit de masse
explique pourquoi l'interaction nucléaire forte a une portée finie et pourquoi les gluons,
bien que sans masse dans la théorie classique, ne produisent que des états liés massifs ---
les simulations sur réseau montrent clairement le déficit, mais une démonstration exige la
limite continue que personne ne sait contrôler. Ce problème demande aux mathématiques de
rattraper la physique.

\begin{refs}
\item Énoncé officiel : Clay Mathematics Institute --- Yang--Mills and the Mass Gap (\url{https://www.claymath.org/millennium/yang-mills-the-maths-gap/})
\item Contexte : Existence de Yang--Mills et déficit de masse --- Wikipédia (en anglais) (\url{https://en.wikipedia.org/wiki/Yang%E2%80%93Mills_existence_and_mass_gap})
\item Lecture : Gerald Folland, \textit{Quantum Field Theory: A Tourist Guide for Mathematicians}
\item Voir aussi : Physique mathématique (\url{applied-physics.html})
\end{refs}

\bigskip

\begin{problem}[{La conjecture de Hodge --- Recherche}]
\textbf{OPEN-05.} \textbf{Problème.} Soit $ X $ une variété projective complexe non
singulière. \textbf{Démontrez que toute classe de Hodge sur $ X $ --- toute classe de
cohomologie rationnelle de type $ (p,p) $ --- est une combinaison linéaire rationnelle de
classes de cohomologie de sous-variétés algébriques complexes de $ X $.} Statut :
ouvert. (La pleine précision demande la théorie de Hodge ; l'énoncé ci-dessus est
l'énoncé standard, et la description du Clay Institute développe la machinerie nécessaire
pour le lire.)
\end{problem}

\noindent William Hodge a proposé cette conjecture au Congrès international de
1950. Elle demande si une condition purement topologico-analytique force l'existence de
géométrie algébrique : étant donné une classe de cohomologie qui \textit{ressemble}
seulement à ce qui pourrait provenir d'une sous-variété, faut-il qu'une sous-variété soit
effectivement là ? La conjecture est connue en degré un, où elle est le théorème de
Lefschetz sur les classes $ (1,1) $, démontré en 1924. Au-delà, on ne sait
essentiellement rien de général, et c'est le moins accessible des sept problèmes --- même
comprendre l'énoncé demande une véritable formation de niveau master, ce qui explique que
les variantes ci-dessous commencent par le théorème de Lefschetz.

\begin{refs}
\item Énoncé officiel : Clay Mathematics Institute --- Hodge Conjecture (\url{https://www.claymath.org/millennium/hodge-conjecture/})
\item Contexte : Conjecture de Hodge --- Wikipédia (\url{https://fr.wikipedia.org/wiki/Conjecture_de_Hodge})
\item Lecture : Claire Voisin, \textit{Théorie de Hodge et géométrie algébrique complexe} (2 vol.)
\end{refs}

\bigskip

\begin{problem}[{La conjecture de Birch et Swinnerton-Dyer --- Recherche}]
\textbf{OPEN-06.} \textbf{Problème.} Soit $ E $ une courbe elliptique sur
$ \mathbb{Q} $, de fonction $ L $ notée $ L(E,s) $. \textbf{Démontrez que le rang du
groupe abélien de type fini $ E(\mathbb{Q}) $ est égal à l'ordre d'annulation de
$ L(E,s) $ en $ s = 1 $.} Statut : ouvert. Connu lorsque l'ordre d'annulation vaut
$ 0 $ ou $ 1 $, grâce aux travaux de Gross--Zagier et de Kolyvagin.
\end{problem}

\noindent Au début des années 1960, Bryan Birch et Peter Swinnerton-Dyer ont
lancé des calculs sur l'ordinateur EDSAC à Cambridge, comptant les points de courbes
elliptiques modulo de nombreux nombres premiers, et ont remarqué un motif : les courbes
ayant beaucoup de solutions rationnelles semblaient avoir en moyenne beaucoup de points
modulo $ p $. La conjecture rend cette observation exacte, en reliant l'arithmétique
d'une courbe au comportement analytique d'une fonction construite à partir d'elle. C'est
l'instance connue la plus profonde d'une philosophie générale --- les fonctions $ L $
encodent l'arithmétique --- et elle a un visage concret : elle réglerait l'antique problème
des nombres congruents, qui demande quels entiers sont l'aire d'un triangle rectangle à
côtés rationnels.

\begin{refs}
\item Énoncé officiel : Clay Mathematics Institute --- Birch and Swinnerton-Dyer Conjecture (\url{https://www.claymath.org/millennium/birch-and-swinnerton-dyer-conjecture/})
\item Contexte : Conjecture de Birch et Swinnerton-Dyer --- Wikipédia (\url{https://fr.wikipedia.org/wiki/Conjecture_de_Birch_et_Swinnerton-Dyer})
\item Lecture : Joseph Silverman \& John Tate, \textit{Rational Points on Elliptic Curves}
\item Voir aussi : Théorie des nombres (\url{pure-number-theory.html})
\end{refs}

\bigskip

\begin{problem}[{La conjecture de Poincaré --- \textit{résolue} --- Recherche}]
\textbf{OPEN-07.} \textbf{Problème (résolu).} \textbf{Démontrez que toute variété close et simplement connexe de
dimension $ 3 $ est homéomorphe à la sphère $ S^3 $ de dimension $ 3 $.} Statut :
\textbf{démontré} par Grigori Perelman dans des prépublications déposées en 2002--2003,
vérifiées par la communauté au cours des trois années suivantes.
\begin{enumerate}
\item[(a)] Lisez la démonstration, et comprenez pourquoi le flot de Ricci avec chirurgie était le
bon outil.
\item[(b)] Attaquez ensuite le cas ouvert : la conjecture de Poincaré lisse en dimension $ 4 $
(voir Topologie (\url{pure-topology.html})).
\end{enumerate}

\end{problem}

\noindent Henri Poincaré a posé la question en 1904, à la naissance même de la
discipline qu'il était en train d'inventer. Les analogues en dimension supérieure sont
tombés les premiers --- Stephen Smale pour la dimension $ \geq 5 $ en 1961, Michael
Freedman pour la dimension 4 (topologiquement) en 1982 ---, laissant le cas original, en
dimension trois, comme le dernier et le plus difficile. Richard Hamilton a proposé le flot
de Ricci, qui lisse la géométrie d'une variété comme l'équation de la chaleur lisse la
température, mais le flot peut former des singularités ; l'idée de Perelman fut de les
contrôler par chirurgie, en découpant et en recollant la variété au fil de son évolution.
Il a ensuite refusé à la fois la médaille Fields (2006) et le million de dollars (2010),
au motif que la récompense était injuste envers Hamilton, dont il jugeait la contribution
non moindre que la sienne. C'est à ce jour le seul problème du millénaire résolu.

\begin{refs}
\item Énoncé officiel : Clay Mathematics Institute --- Poincaré Conjecture (\url{https://www.claymath.org/millennium/poincare-conjecture/})
\item Contexte : Conjecture de Poincaré --- Wikipédia (\url{https://fr.wikipedia.org/wiki/Conjecture_de_Poincar%C3%A9})
\item Article : Grigori Perelman, \og{} The entropy formula for the Ricci flow and its geometric applications \fg{} (2002), arXiv:math/0211159 (\url{https://arxiv.org/abs/math/0211159})
\item Lecture : Donal O'Shea, \textit{The Poincaré Conjecture} (2007)
\end{refs}

\bigskip

\begin{problem}[{Les trois premiers zéros, à partir d'une table --- Lycée, short}]
\textbf{OPEN-43.} \textbf{Problème.} Dans une table publiée des zéros de la fonction zêta
de Riemann (Andrew Odlyzko tient en ligne les tables de référence), relevez les parties
imaginaires des trois premiers zéros non triviaux --- vous devriez trouver environ
$ 14{,}1347 $, $ 21{,}0220 $ et $ 25{,}0109 $ --- et vérifiez que la table donne
chaque zéro avec une partie réelle exactement égale à $ \tfrac{1}{2} $. Expliquez
ensuite en une phrase pourquoi $ 10^{13} $ confirmations de ce genre ne démontreraient
toujours pas l'hypothèse de Riemann.
\end{problem}

\noindent Savoir où vivent les données est une vraie compétence. L'étude
numérique des zéros de zêta va des calculs à la main de Riemann lui-même à Alan Turing,
qui a calculé des zéros sur le Manchester Mark 1 en 1950, jusqu'aux immenses tables
modernes d'Odlyzko. Manipuler trois nombres réels rend OPEN-01 concret comme aucun résumé
ne saurait le faire.

\begin{refs}
\item Contexte : Hypothèse de Riemann --- Wikipédia (\url{https://fr.wikipedia.org/wiki/Hypoth%C3%A8se_de_Riemann})
\item Tables : Andrew Odlyzko --- Tables of zeros of the Riemann zeta function (\url{https://www-users.cse.umn.edu/\textasciitilde{}odlyzko/zeta_tables/index.html})
\end{refs}

\bigskip

\begin{problem}[{P, NP et co-NP en trois phrases --- Licence, short}]
\textbf{OPEN-44.} \textbf{Problème.} Énoncez la différence entre $ \mathrm{P} $,
$ \mathrm{NP} $ et co-NP en trois phrases --- une par classe, formulées en termes de
certificats \og{} oui \fg{} et de certificats \og{} non \fg{}. Nommez ensuite un problème de décision
naturel qui appartient à $ \mathrm{NP} \cap \text{co-}\mathrm{NP} $ sans qu'on sache
s'il appartient à $ \mathrm{P} $.
\end{problem}

\noindent Savoir énoncer proprement ces définitions est le droit d'entrée pour
OPEN-02. L'exemple standard est la factorisation, posée comme problème de décision :
exhiber un facteur pour \og{} oui \fg{}, exhiber la décomposition complète en facteurs premiers
pour \og{} non \fg{} --- la primalité elle-même est dans $ \mathrm{P} $ depuis
Agrawal--Kayal--Saxena (2002). Savoir si $ \mathrm{NP} = \text{co-}\mathrm{NP} $ est
lui-même ouvert, et une démonstration que ces classes diffèrent donnerait déjà
$ \mathrm{P} \neq \mathrm{NP} $.

\begin{refs}
\item Contexte : co-NP --- Wikipédia (\url{https://fr.wikipedia.org/wiki/Co-NP})
\item Lecture : Michael Sipser, \textit{Introduction to the Theory of Computation}, chap. 7
\end{refs}

\bigskip

\begin{problem}[{Le changement d'échelle qui rend Navier--Stokes difficile --- Master, short}]
\textbf{OPEN-45.} \textbf{Problème.} Montrez que si le couple $ (u, p) $ est solution
des équations de Navier--Stokes tridimensionnelles d'OPEN-03, alors, pour tout
$ \lambda > 0 $, il en va de même du couple redimensionné
\[ u_{\lambda}(x,t) = \lambda\, u(\lambda x, \lambda^{2} t), \qquad
p_{\lambda}(x,t) = \lambda^{2} p(\lambda x, \lambda^{2} t) . \]
Déterminez ensuite l'exposant $ s $ pour lequel la norme de la donnée initiale
$ \| u(\cdot,0) \|_{\dot{H}^{s}(\mathbb{R}^{3})} $ est invariante par ce changement
d'échelle, et expliquez en une phrase pourquoi la norme d'énergie $ L^{2} $, qui est
conservée, est de ce fait dite \textit{surcritique} en dimension trois.
\end{problem}

\noindent Ce calcul d'une demi-page explique toute la forme du problème du
millénaire. La seule quantité dont on sache qu'elle est contrôlée pour tout temps est
l'énergie, et le changement d'échelle ci-dessus montre que la norme d'énergie se situe du
mauvais côté de l'exposant critique en dimension trois --- si bien que rétrécir une solution
la fait paraître, du point de vue de l'énergie, de plus en plus petite tandis que ses
gradients explosent sans entrave. En dimension deux, le même calcul place l'énergie
exactement au niveau critique, ce qui explique que le théorème de Ladyjenskaïa existe et
que son analogue tridimensionnel n'existe pas. La construction d'explosion moyennée de
Terence Tao est l'énoncé le plus net de la difficulté : toute démonstration de régularité
globale devra faire appel à autre chose que l'énergie et le changement d'échelle.

\begin{refs}
\item Contexte : Solutions des équations de Navier-Stokes --- Wikipédia (\url{https://fr.wikipedia.org/wiki/Solutions_des_%C3%A9quations_de_Navier-Stokes})
\item Lecture : Terence Tao, \og{} Finite time blowup for an averaged three-dimensional Navier-Stokes equation \fg{}, \textit{Journal of the AMS} 29 (2016), arXiv:1402.0290 (\url{https://arxiv.org/abs/1402.0290})
\item Voir aussi : Équations aux dérivées partielles (\url{applied-pde.html})
\end{refs}

\bigskip


\section*{Des variantes réellement abordables}

\begin{problem}[{Le produit d'Euler, ou pourquoi $ \zeta $ connaît les nombres premiers --- Licence}]
\textbf{OPEN-08.} \textbf{Problème.} Soit $ s > 1 $ un réel.
\begin{enumerate}
\item[(a)] Démontrez la formule du produit d'Euler :
\[ \sum_{n=1}^{\infty} \frac{1}{n^s} = \prod_{p \text{ premier}} \frac{1}{1 - p^{-s}} . \]
\item[(b)] Déduisez-en, en faisant tendre $ s \to 1^{+} $, que $ \sum_{p} 1/p $ diverge.
\item[(c)] Expliquez pourquoi cela fournit une seconde démonstration de l'infinité des nombres
premiers --- plus forte que celle d'Euclide, puisqu'elle dit que les nombres premiers ne
sont pas trop rares.
\end{enumerate}

\end{problem}

\noindent Euler a trouvé cette formule en 1737, et c'est la seule formule dont
naisse toute la théorie analytique des nombres : le membre de gauche est de l'analyse,
celui de droite de l'arithmétique, et l'égalité est le théorème fondamental de
l'arithmétique déguisé. Tout ce qui concerne l'hypothèse de Riemann commence ici.
Comprendre exactement pourquoi l'unicité de la factorisation produit un produit sur les
nombres premiers est le prérequis pour comprendre ce qu'un zéro de $ \zeta $ peut bien
avoir à faire avec la répartition des nombres premiers.

\begin{refs}
\item Contexte : Produit eulérien --- Wikipédia (\url{https://fr.wikipedia.org/wiki/Produit_eul%C3%A9rien})
\item Lecture : Tom Apostol, \textit{Introduction to Analytic Number Theory}, chap. 11
\item Voir aussi : Théorie des nombres (\url{pure-number-theory.html}), L'art de la démonstration (\url{proofs.html})
\end{refs}

\bigskip

\begin{problem}[{Aucun zéro sur la droite $ \operatorname{Re}(s) = 1 $ --- Master}]
\textbf{OPEN-09.} \textbf{Problème.}
\begin{enumerate}
\item[(a)] Travaillez et reconstituez la démonstration du fait que $ \zeta(1 + it) \neq 0 $
pour tout réel $ t \neq 0 $, en utilisant l'inégalité
\[ 3 + 4\cos\theta + \cos 2\theta \geq 0 . \]
\item[(b)] Expliquez précisément comment cette non-annulation donne le théorème des nombres
premiers, $ \pi(x) \sim x / \log x $.
\item[(c)] Énoncez ce qu'améliorerait l'hypothèse de Riemann : donnez le terme d'erreur qu'elle
implique pour $ \pi(x) $.
\end{enumerate}

\end{problem}

\noindent C'est le théorème de Hadamard--de la Vallée Poussin de 1896, le résultat
qui a fait de la théorie analytique des nombres une discipline. L'astuce
trigonométrique ressemble à un lapin sorti d'un chapeau ; il vaut la peine de rester avec
elle jusqu'à ce qu'elle cesse d'y ressembler. Faire cet exercice enseigne la véritable
leçon de l'hypothèse de Riemann : le théorème que vous savez démontrer écarte les zéros
d'une droite, la conjecture que vous ne savez pas démontrer les place tous sur une autre,
et la distance entre ces deux énoncés est la distance entre savoir que les nombres
premiers sont réguliers et savoir exactement à quel point.

\begin{refs}
\item Contexte : Théorème des nombres premiers --- Wikipédia (\url{https://fr.wikipedia.org/wiki/Th%C3%A9or%C3%A8me_des_nombres_premiers})
\item Lecture : Tom Apostol, \textit{Introduction to Analytic Number Theory}, chap. 13
\item Lecture : Harold Davenport, \textit{Multiplicative Number Theory}
\end{refs}

\bigskip

\begin{problem}[{2-SAT est facile, 3-SAT ne l'est pas --- Licence}]
\textbf{OPEN-10.} \textbf{Problème.} Rappelons qu'une formule sous forme normale conjonctive est une
conjonction de clauses, chacune étant une disjonction de littéraux.
\begin{enumerate}
\item[(a)] Démontrez que $ \mathrm{P} \subseteq \mathrm{NP} $.
\item[(b)] Donnez un algorithme en temps polynomial qui décide la satisfiabilité d'une formule
booléenne sous forme normale conjonctive comportant exactement deux littéraux par clause,
et démontrez qu'il est correct --- la méthode du graphe d'implications, où une formule est
insatisfiable exactement lorsqu'une variable $ x $ et sa négation appartiennent à la
même composante fortement connexe.
\item[(c)] Énoncez le théorème de Cook--Levin et expliquez, sans le démontrer intégralement,
pourquoi le passage de 3-SAT dans $ \mathrm{P} $ ferait s'effondrer tout
$ \mathrm{NP} $.
\end{enumerate}

\end{problem}

\noindent L'écart entre deux littéraux et trois, c'est tout le mystère de
$ \mathrm{P} $ contre $ \mathrm{NP} $ en miniature : l'un se résout en temps linéaire,
l'autre est le problème difficile canonique, et rien dans leurs énoncés n'explique la
différence. Cook a démontré en 1971 que SAT est NP-complet ; l'article de Karp de 1972 a
aussitôt produit vingt et un autres problèmes NP-complets et déclenché le déluge. Sentir
exactement où la difficulté s'introduit quand on ajoute le troisième littéral, c'est le
commencement de l'intuition en complexité.

\begin{refs}
\item Contexte : Problème 2-SAT --- Wikipédia (\url{https://fr.wikipedia.org/wiki/Probl%C3%A8me_2-SAT})
\item Contexte : Théorème de Cook --- Wikipédia (\url{https://fr.wikipedia.org/wiki/Th%C3%A9or%C3%A8me_de_Cook})
\item Lecture : Michael Sipser, \textit{Introduction to the Theory of Computation}, chap. 7
\item Cours : MIT OCW 6.045J Automata, Computability, and Complexity (\url{https://ocw.mit.edu/courses/6-045j-automata-computability-and-complexity-spring-2011/})
\end{refs}

\bigskip

\begin{problem}[{La barrière de la relativisation --- Master}]
\textbf{OPEN-11.} \textbf{Problème.} Une machine à oracle peut interroger l'appartenance à un ensemble fixé
pour un coût unitaire ; $ \mathrm{P}^{A} $ et $ \mathrm{NP}^{A} $ désignent les
classes relativisées à l'oracle $ A $.
\begin{enumerate}
\item[(a)] Reconstituez le théorème de Baker--Gill--Solovay : construisez des oracles $ A $ et
$ B $ tels que $ \mathrm{P}^{A} = \mathrm{NP}^{A} $ et
$ \mathrm{P}^{B} \neq \mathrm{NP}^{B} $.
\item[(b)] Expliquez la conséquence : pourquoi toute technique de démonstration qui \og{} relativise \fg{}
--- qui reste valable lorsqu'on remet le même oracle aux deux machines --- est incapable de
trancher $ \mathrm{P} $ contre $ \mathrm{NP} $, et pourquoi la diagonalisation seule
ne peut donc pas régler la question.
\end{enumerate}

\end{problem}

\noindent Ce résultat de 1975 explique pourquoi la théorie de la complexité est
difficile d'une manière qui est elle-même démontrable. Il ne dit pas que le problème est
insoluble ; il dit qu'une famille naturelle entière d'attaques est morte, et il a tué
l'optimisme selon lequel les arguments de simulation et de diagonalisation, qui avaient si
bien fonctionné pour le problème de l'arrêt et pour le théorème de hiérarchie en temps,
règleraient tout. Razborov et Rudich ont plus tard ajouté une deuxième barrière (les
démonstrations naturelles, 1994), et Aaronson et Wigderson une troisième
(l'algébrisation, 2008). Savoir ce qui ne peut pas marcher constitue l'essentiel du savoir
d'un chercheur.

\begin{refs}
\item Contexte : Oracle (machine de Turing) --- Wikipédia (\url{https://fr.wikipedia.org/wiki/Oracle_%28machine_de_Turing%29})
\item Lecture : Arora \& Barak, \textit{Computational Complexity: A Modern Approach}, chap. 3
\end{refs}

\bigskip

\begin{problem}[{La dimension deux se porte bien : les solutions faibles de Leray --- Master}]
\textbf{OPEN-12.} \textbf{Problème.} Considérez les équations de Navier--Stokes d'OPEN-03.
\begin{enumerate}
\item[(a)] Définissez une solution faible de Leray--Hopf des équations de Navier--Stokes et
démontrez l'inégalité d'énergie qu'elle vérifie.
\item[(b)] Lisez et reconstituez la démonstration du fait qu'en dimension \textit{deux} les
solutions faibles sont uniques et régulières pour tout temps.
\item[(c)] Identifiez, précisément et en un paragraphe, l'étape de cet argument qui échoue en
dimension trois --- le terme d'étirement de tourbillon $ (\omega \cdot \nabla) u $, qui
est identiquement nul en dimension 2.
\end{enumerate}

\end{problem}

\noindent Jean Leray a construit les solutions faibles en 1934, en inventant pour
cela une bonne part de la théorie moderne des EDP, et le fait que le problème
bidimensionnel soit complètement compris tandis que le tridimensionnel reste un mystère à
un million de dollars est ce qu'il y a de plus instructif dans tout le sujet. La question (c)
est le véritable exercice : localiser la phrase exacte où une démonstration cesse de
fonctionner apprend davantage que n'importe quel nombre de démonstrations réussies. Toute
attaque sérieuse du problème du millénaire est une tentative de contrôler ce terme unique.

\begin{refs}
\item Contexte : Équations de Navier-Stokes --- Wikipédia (\url{https://fr.wikipedia.org/wiki/%C3%89quations_de_Navier-Stokes})
\item Lecture : Lawrence C. Evans, \textit{Partial Differential Equations}, chap. 7
\item Voir aussi : Équations aux dérivées partielles (\url{applied-pde.html})
\end{refs}

\bigskip

\begin{problem}[{Des chocs en dimension un : l'équation de Burgers --- Licence}]
\textbf{OPEN-13.} \textbf{Problème.} Considérez l'équation de Burgers sans viscosité
\[ \partial_t u + u \, \partial_x u = 0, \qquad u(x,0) = u_0(x). \]
\begin{enumerate}
\item[(a)] Utilisez la méthode des caractéristiques pour montrer que si $ u_0 $ est lisse et
strictement décroissante quelque part, la solution développe une discontinuité en temps
fini, et calculez exactement cet instant en fonction de $ \min u_0' $.
\item[(b)] Expliquez pourquoi cela ne règle \textit{pas} la question de Navier--Stokes --- que fait
le terme visqueux $ \nu \Delta u $ qui change la donne ?
\end{enumerate}

\end{problem}

\noindent C'est la rencontre honnête la moins coûteuse avec l'explosion en temps
fini : une onde parfaitement lisse qui se redresse jusqu'à déferler, exactement comme les
vagues de l'océan. Cela montre que le transport non linéaire à lui seul détruit la
régularité, et cela rend vivante la vraie question du problème du millénaire --- la
viscosité gagne-t-elle toujours la course contre la non-linéarité en dimension trois ?
Mener ce calcul à la main, c'est toute la différence entre connaître l'expression
\og{} explosion en temps fini \fg{} et la comprendre.

\begin{refs}
\item Contexte : Équation de Burgers --- Wikipédia (\url{https://fr.wikipedia.org/wiki/%C3%89quation_de_Burgers})
\item Lecture : Walter Strauss, \textit{Partial Differential Equations: An Introduction}, chap. 14
\end{refs}

\bigskip

\begin{problem}[{Théorie de jauge sur réseau : la version finie de Yang--Mills --- Master}]
\textbf{OPEN-14.} \textbf{Problème.} Définissez la théorie de jauge sur réseau de Wilson pour un groupe
compact $ G $ sur un réseau périodique fini : le champ est une attribution d'éléments du
groupe aux arêtes, et l'action est une somme sur les plaquettes.
\begin{enumerate}
\item[(a)] Démontrez que la mesure obtenue est bien définie (c'est une intégrale en dimension
finie --- aucune difficulté analytique).
\item[(b)] Démontrez le confinement dans la limite de couplage fort : montrez que l'espérance de
la boucle de Wilson décroît comme l'exponentielle de l'aire enclose.
\item[(c)] Énoncez précisément ce qu'exigerait la limite continue, et pourquoi personne ne sait
la prendre.
\end{enumerate}

\end{problem}

\noindent Kenneth Wilson a introduit le réseau en 1974, et c'est ainsi que
l'interaction forte est effectivement calculée aujourd'hui, sur des superordinateurs. Le
fait pédagogique frappant est que la théorie sur réseau est \textit{facile} --- un objet fini,
honnête et rigoureux, qu'un étudiant de licence peut définir --- et que tout le problème du
millénaire tient dans l'unique étape consistant à faire tendre le pas du réseau vers zéro
en gardant la physique fixe. Ce problème vous montre exactement où se cache le million de
dollars.

\begin{refs}
\item Contexte : Théorie de jauge sur réseau --- Wikipédia (\url{https://fr.wikipedia.org/wiki/Th%C3%A9orie_de_jauge_sur_r%C3%A9seau})
\item Contexte : Boucle de Wilson --- Wikipédia (\url{https://fr.wikipedia.org/wiki/Boucle_de_Wilson})
\item Lecture : Barry Simon, \textit{The Statistical Mechanics of Lattice Gases}
\end{refs}

\bigskip

\begin{problem}[{Le théorème de Lefschetz $ (1,1) $ : Hodge en degré un --- Master}]
\textbf{OPEN-15.} \textbf{Problème.} Soit $ X $ une variété kählérienne compacte.
\begin{enumerate}
\item[(a)] Lisez et reconstituez la démonstration du théorème de Lefschetz $ (1,1) $ : une
classe de $ H^2(X, \mathbb{Z}) $ est la première classe de Chern d'un fibré en droites
holomorphe si et seulement si son image dans $ H^2(X, \mathbb{C}) $ est de type
$ (1,1) $.
\item[(b)] Énoncez la conjecture de Hodge et vérifiez que ce théorème en est exactement le cas
de degré un.
\end{enumerate}

\end{problem}

\noindent Démontré par Solomon Lefschetz en 1924, trois décennies avant que Hodge
ne pose sa conjecture, c'est l'unique cas où le motif est confirmé --- et sa démonstration,
via la suite exponentielle de faisceaux, est un morceau de mathématiques beau et complet
qui enseigne en outre le langage dans lequel la conjecture générale est écrite. Si vous
voulez comprendre ce que la conjecture de Hodge affirme seulement, voici la porte.

\begin{refs}
\item Contexte : Théorème de Lefschetz sur les classes (1,1) --- Wikipédia (en anglais) (\url{https://en.wikipedia.org/wiki/Lefschetz_theorem_on_(1,1)-classes})
\item Lecture : Phillip Griffiths \& Joseph Harris, \textit{Principles of Algebraic Geometry}, chap. 1
\item Lecture : Claire Voisin, \textit{Théorie de Hodge et géométrie algébrique complexe I}
\end{refs}

\bigskip

\begin{problem}[{Les nombres congruents : le visage concret de BSD --- Licence}]
\textbf{OPEN-16.} \textbf{Problème.} Un entier strictement positif $ n $ est dit \textit{congruent} s'il
est l'aire d'un triangle rectangle à côtés rationnels.
\begin{enumerate}
\item[(a)] Montrez que $ 6 $ est congruent, et montrez que $ 5 $ l'est en exhibant un
triangle explicite.
\item[(b)] Démontrez que $ n $ est congruent si et seulement si la courbe elliptique
$ y^2 = x^3 - n^2 x $ possède un point rationnel d'ordre infini.
\item[(c)] Utilisez le théorème de Tunnell (1983) pour tester si $ n = 157 $ est congruent, en
comptant les solutions de deux formes quadratiques ternaires --- un calcul fini. Notez
honnêtement ce que vous avez fait : le critère de Tunnell est démontré inconditionnellement
dans un sens, et sa réciproque dépend de la conjecture de Birch et Swinnerton-Dyer.
\end{enumerate}

\end{problem}

\noindent Ce problème a mille ans --- il apparaît dans un manuscrit arabe des
environs de l'an 900 et Fibonacci en a discuté en 1225 --- et il n'est toujours pas
entièrement résolu, parce que sa solution complète équivaut à un cas d'un problème du
millénaire. La question (b) est le moment où l'abstraction paie : une question antique sur
des triangles devient une question sur les points rationnels d'une courbe, et voilà que
toute la machinerie de la géométrie arithmétique s'applique. Le calcul de Don Zagier
montrant que le plus simple triangle rectangle d'aire $ 157 $ a une hypoténuse dont le
numérateur compte 45 chiffres est la démonstration classique que \og{} il n'y a qu'à
chercher \fg{} n'avait aucune chance de marcher.

\begin{refs}
\item Contexte : Nombre congruent --- Wikipédia (\url{https://fr.wikipedia.org/wiki/Nombre_congruent})
\item Contexte : Théorème de Tunnell --- Wikipédia (en anglais) (\url{https://en.wikipedia.org/wiki/Tunnell%27s_theorem})
\item Lecture : Neal Koblitz, \textit{Introduction to Elliptic Curves and Modular Forms}, chap. 1
\item Voir aussi : Théorie des nombres (\url{pure-number-theory.html})
\end{refs}

\bigskip

\begin{problem}[{Un nombre composé qui passe le test de Fermat --- Licence, short}]
\textbf{OPEN-46.} \textbf{Problème.} Vérifiez à la main que
\[ 2^{340} \equiv 1 \pmod{341}, \qquad 341 = 11 \times 31 , \]
le chemin le plus rapide étant de remarquer que $ 2^{10} = 1024 = 3 \cdot 341 + 1 $.
Énoncez ensuite précisément ce que cet unique exemple démontre au sujet de la réciproque
du petit théorème de Fermat, et vérifiez que la base $ 3 $ ne se laisse pas tromper en
calculant $ 3^{340} \bmod 341 $.
\end{problem}

\noindent Le petit théorème de Fermat dit que $ a^{p-1} \equiv 1 \pmod p $ pour
$ p $ premier ne divisant pas $ a $ ; la réciproque tentante est fausse, et $ 341 $
--- trouvé par Pierre Frédéric Sarrus en 1819 --- en est le plus petit témoin pour la base
$ 2 $. Pire, les nombres de Carmichael comme $ 561 = 3 \cdot 11 \cdot 17 $ passent le
test pour toute base première avec eux, raison pour laquelle les vrais tests de primalité
utilisent le raffinement de Miller--Rabin. C'est dans ce raffinement que se rencontrent les
deux plus grands problèmes de cette page : Miller a montré en 1976 qu'une version
déterministe en temps polynomial découle de l'hypothèse de Riemann généralisée (OPEN-01),
et il a fallu attendre Agrawal--Kayal--Saxena en 2002 pour obtenir un test inconditionnel en
temps polynomial (OPEN-02).

\begin{refs}
\item Contexte : Nombre pseudo-premier --- Wikipédia (\url{https://fr.wikipedia.org/wiki/Nombre_pseudo-premier})
\item Contexte : Test de primalité de Miller-Rabin --- Wikipédia (\url{https://fr.wikipedia.org/wiki/Test_de_primalit%C3%A9_de_Miller-Rabin})
\item Lecture : Richard Crandall \& Carl Pomerance, \textit{Prime Numbers: A Computational Perspective}, chap. 3
\item Voir aussi : Cryptographie (\url{applied-crypto.html})
\end{refs}

\bigskip

\begin{problem}[{La constante de Legendre, ou ce qu'un ajustement de courbe ne peut pas voir --- Licence}]
\textbf{OPEN-57.} \textbf{Problème.} En 1808, Adrien-Marie Legendre, ajustant une formule à des tables de
nombres premiers, a proposé l'approximation
\[ \pi(x) \;\approx\; \frac{x}{\log x - B}, \qquad B = 1.08366 . \]
Le nombre B s'appelle la constante de Legendre.
\begin{enumerate}
\item[(a)] Comparez les deux candidates $ x/(\log x - 1.08366) $ et $ x/(\log x - 1) $ aux
valeurs exactes de $ \pi(x) $ en $ x = 10^3, 10^4, 10^5, 10^6 $ (elles sont tabulées
partout), et dites laquelle s'ajuste le mieux sur cet intervalle.
\item[(b)] Démontrez l'observation de Tchebychev, en 1849 : si la limite
\[ B \;=\; \lim_{x \to \infty} \left( \log x - \frac{x}{\pi(x)} \right) \]
existe, alors elle vaut 1. Montrez ensuite que le théorème des nombres premiers, sous la
forme fine $ \pi(x) \sim \operatorname{li}(x) $, force cette limite à exister.
\item[(c)] Expliquez pourquoi Legendre n'a pas seulement manqué de chance. Montrez que les deux
formules candidates diffèrent d'une quantité négligeable devant $ \pi(x) $ quand
$ x \to \infty $, et pourtant grande en valeur absolue sur tout l'intervalle où Legendre
pouvait calculer, puis tirez-en la conclusion méthodologique sur ce qu'un ajustement
numérique peut et ne peut pas déterminer.
\item[(d)] Comparez avec la conjecture de Gauss. Démontrez que
\[ \operatorname{li}(x) \;=\; \frac{x}{\log x} + \frac{x}{(\log x)^2} + O\!\left( \frac{x}{(\log x)^3} \right), \]
et expliquez pourquoi le logarithme intégral, qui ne comporte aucun paramètre ajustable,
bat la formule de Legendre, qui en comporte un.
\end{enumerate}

\end{problem}

\noindent Legendre a publié son approximation dans l'\textit{Essai sur la théorie des
nombres} et a affiné la constante à 1,08366 dans l'édition de 1808 ; Gauss avait deviné
en privé le logarithme intégral, adolescent, vers 1792, et n'avait rien publié --- de sorte
que les deux meilleures conjectures du dix-huitième siècle ont été formées indépendamment
et qu'une seule fut défendue par écrit. Tchebychev a démontré en 1849 qu'aucune constante
autre que 1 n'était possible, et de la Vallée Poussin a réglé la question en 1899, comme
sous-produit de son terme d'erreur dans le théorème des nombres premiers. La valeur
1,08366 est un fossile de surajustement : c'est le nombre qui colle le mieux aux données
dont Legendre disposait, et il se trompe sur l'univers. C'est exactement pour cela que la
théorie analytique des nombres énonce ses résultats de façon asymptotique et se méfie des
tables --- une discipline dont dépend tout le reste de cette page, et qu'il vaut la peine de
garder en mémoire à côté de la démonstration par Littlewood en 1914 que
$ \operatorname{li}(x) - \pi(x) $ change de signe une infinité de fois, alors qu'il est
positif pour tous les x que quiconque a jamais calculés.

\begin{refs}
\item Contexte : Constante de Legendre --- Wikipédia (\url{https://fr.wikipedia.org/wiki/Constante_de_Legendre})
\item Contexte : Théorème des nombres premiers --- Wikipédia (\url{https://fr.wikipedia.org/wiki/Th%C3%A9or%C3%A8me_des_nombres_premiers})
\item Lecture : Tom Apostol, \textit{Introduction to Analytic Number Theory}, chap. 4
\item Voir aussi : Théorie des nombres (\url{pure-number-theory.html})
\end{refs}

\bigskip

\begin{problem}[{Une hypothèse de Riemann qui est un théorème --- Master}]
\textbf{OPEN-58.} \textbf{Problème.} Soit C une courbe projective lisse de genre g sur le corps fini
$ \mathbb{F}_q $, et définissons sa fonction zêta par
\[ Z_C(T) \;=\; \exp\left( \sum_{n \ge 1} \frac{\# C(\mathbb{F}_{q^n})}{n} T^n \right). \]
Dans ce cadre, l'analogue d'OPEN-01 n'est pas une conjecture. C'est un théorème, et
l'objet de ce problème est de découvrir pourquoi.
\begin{enumerate}
\item[(a)] Calculez explicitement $ Z_C(T) $ pour la droite projective
$ C = \mathbb{P}^1 $, et vérifiez à la main que c'est une fraction rationnelle en T
satisfaisant l'équation fonctionnelle attendue reliant T et $ 1/(qT) $.
\item[(b)] Énoncez les conjectures de Weil pour les courbes --- rationalité, équation fonctionnelle
et l'\og{} hypothèse de Riemann \fg{} selon laquelle tout zéro de $ Z_C $ a pour module
$ q^{-1/2} $ --- et démontrez que la dernière équivaut à la borne de Hasse--Weil
\[ \bigl| \# C(\mathbb{F}_q) - (q+1) \bigr| \;\le\; 2g\sqrt{q}. \]
\item[(c)] Démontrez le cas du genre un, qui est le théorème de Hasse de 1933 : pour une courbe
elliptique E sur $ \mathbb{F}_q $,
$ |\# E(\mathbb{F}_q) - (q+1)| \le 2\sqrt{q} $. (Chemin : l'endomorphisme de Frobenius
vérifie une équation du second degré, et l'application degré sur les endomorphismes est
une forme quadratique définie positive ; appliquez ensuite l'inégalité de
Cauchy--Schwarz.)
\item[(d)] Dites maintenant ce qui manque sur $\mathbb{Q}$. Nommez l'objet qui joue le rôle du Frobenius
dans le tableau des corps de fonctions, identifiez l'objet géométrique sur lequel il agit,
et expliquez en un paragraphe pourquoi ni l'un ni l'autre n'a de contrepartie connue pour
la fonction zêta de Riemann ordinaire. Voilà la réponse honnête à \og{} pourquoi OPEN-01
est-il toujours ouvert \fg{}.
\end{enumerate}

\end{problem}

\noindent Voici la chose la plus utile à savoir sur l'hypothèse de Riemann : dans
un univers parallèle, c'est un théorème, et la démonstration est comprise. Emil Artin a
remarqué l'analogie dans sa thèse de 1921 en attachant des fonctions zêta aux corps de
fonctions ; Hasse a démontré le cas du genre un dans les années 1930 ; André Weil a
annoncé le cas général pour les courbes en 1940 --- depuis une cellule de prison à Rouen, où
il était détenu pour ne pas s'être présenté au service militaire --- et a publié des
démonstrations complètes en 1948, inventant au passage une grande partie de la géométrie
algébrique moderne. Grothendieck a ensuite refondé la discipline entière pour attaquer les
conjectures en dimension supérieure, et Deligne les a achevées en 1974. Corps de nombres
et corps de fonctions se ressemblent presque partout ; la différence décisive est qu'une
courbe sur $ \mathbb{F}_q $ vient avec une application de Frobenius et un espace où la
faire agir, tandis que $\mathbb{Q}$ ne vient ni avec l'une ni avec l'autre. La longue quête d'un
\og{} corps à un élément \fg{} et d'une cohomologie à la Weil sur $\mathbb{Z}$ est un programme de recherche
sérieux bâti précisément sur cet écart.

\begin{refs}
\item Contexte : Conjectures de Weil --- Wikipédia (\url{https://fr.wikipedia.org/wiki/Conjectures_de_Weil})
\item Contexte : Théorème de Hasse sur les courbes elliptiques --- Wikipédia (\url{https://fr.wikipedia.org/wiki/Th%C3%A9or%C3%A8me_de_Hasse_sur_les_courbes_elliptiques})
\item Lecture : Joseph H. Silverman, \textit{The Arithmetic of Elliptic Curves}, chap. V
\item Lecture : Kenneth Ireland \& Michael Rosen, \textit{A Classical Introduction to Modern Number Theory}, chap. 11
\end{refs}

\bigskip

\begin{problem}[{Robin et Lagarias : l'hypothèse de Riemann en habits élémentaires --- Licence}]
\textbf{OPEN-59.} \textbf{Problème.} Soit $ \sigma(n) = \sum_{d \mid n} d $ la somme des diviseurs et
soit $ H_n = 1 + \tfrac12 + \cdots + \tfrac1n $. Chacun des deux énoncés élémentaires
suivants est \textit{équivalent} à l'hypothèse de Riemann d'OPEN-01. Robin (1984) :
l'inégalité $ \sigma(n) 5040 $.
Lagarias (2002) : l'inégalité $ \sigma(n) \le H_n + e^{H_n} \log H_n $ vaut pour tout
$ n \ge 1 $, avec égalité seulement en $ n = 1 $.
\begin{enumerate}
\item[(a)] Mettez les mains dans le cambouis. Testez l'inégalité de Robin sur les nombres riches
en diviseurs $ n = 12, 60, 120, 360, 2520, 5040 $ et sur quelques nombres premiers, puis
trouvez tous les $ n \le 5040 $ qui la violent --- l'ensemble exceptionnel est court et
célèbre.
\item[(b)] Démontrez le théorème de Gronwall (1913), ou reconstituez-le à partir des références :
\[ \limsup_{n \to \infty} \frac{\sigma(n)}{n \log \log n} \;=\; e^{\gamma}. \]
Concluez que l'inégalité de Robin affirme exactement que cette borne supérieure est
approchée sans jamais être atteinte au-delà de 5040 --- le critère est donc délicat par
construction, non par accident.
\item[(c)] Montrez qu'il suffit de vérifier l'inégalité de Robin sur les nombres colossalement
abondants, et expliquez la réduction : pourquoi un contre-exemple doit-il nécessairement
battre un record pour $ \sigma(n)/n $ ?
\item[(d)] Répondez ensuite à la question honnête. Déduisez l'inégalité de Lagarias de celle de
Robin pour n grand, ou suivez sa dérivation dans la référence ; puis écrivez un paragraphe
sur ce qui a été gagné, si quelque chose l'a été. Les deux énoncés sont élémentaires ;
tous deux sont équivalents à un problème que personne ne sait résoudre. Que cela vous
apprend-il sur l'expression \og{} formulation équivalente \fg{} ?
\end{enumerate}

\end{problem}

\noindent Guy Robin a démontré son critère en 1984 dans un article sur l'ordre
maximal de la fonction somme des diviseurs --- mais il ne fut pas le premier à trouver
l'inégalité. Ramanujan avait démontré essentiellement la même chose, sous l'hypothèse de
Riemann, dans son mémoire de 1915 sur les nombres hautement composés ; les pages
correspondantes furent coupées de la version publiée à cause de la pénurie de papier due à
la guerre, et ne parurent qu'en 1997, quand Nicolas et Robin publièrent le manuscrit
supprimé. La note de Lagarias parue en 2002 dans le \textit{Monthly} pousse la reformulation
aussi loin qu'elle peut aller : des nombres harmoniques et une somme de diviseurs, pas la
moindre analyse complexe en vue, et pourtant l'un des problèmes ouverts les plus difficiles
des mathématiques. Voilà la leçon à retenir. Une formulation équivalente redistribue la
difficulté, elle ne la dissout pas --- et qu'un énoncé soit élémentaire ne dit absolument
rien sur le caractère élémentaire d'une démonstration.

\begin{refs}
\item Contexte : Fonction somme des diviseurs (théorème de Robin) --- Wikipédia (\url{https://fr.wikipedia.org/wiki/Fonction_somme_des_puissances_k-i%C3%A8mes_des_diviseurs})
\item Contexte : Nombre colossalement abondant --- Wikipédia (\url{https://fr.wikipedia.org/wiki/Nombre_colossalement_abondant})
\item Article : Jeffrey C. Lagarias, \og{} An elementary problem equivalent to the Riemann hypothesis \fg{}, \textit{American Mathematical Monthly} 109 (2002), arXiv:math/0008177 (\url{https://arxiv.org/abs/math/0008177})
\item Article : Guy Robin, \og{} Grandes valeurs de la fonction somme des diviseurs et hypothèse de Riemann \fg{}, \textit{Journal de Mathématiques Pures et Appliquées} 63 (1984)
\end{refs}

\bigskip

\begin{problem}[{Remplacez le 3 par un 5 --- Licence, short}]
\textbf{OPEN-60.} \textbf{Problème.} Définissez $ U(n) = n/2 $ pour n pair et
$ U(n) = 5n + 1 $ pour n impair, puis calculez les trajectoires de $ n = 5 $,
$ n = 7 $ et $ n = 13 $ sur cent pas chacune. Dérivez ensuite le multiplicateur
heuristique par pas impair pour U et pour l'application de Collatz T d'OPEN-17 --- en
moyenne, combien de facteurs 2 contient $ 3n+1 $, et combien en contient $ 5n+1 $ ? ---
et dites laquelle des deux applications vous vous attendriez à voir envoyer presque toute
valeur initiale vers l'infini.
\end{problem}

\noindent Le calcul prend quelques minutes et réorganise durablement votre
intuition. Deux des trois valeurs initiales ci-dessus tombent dans un cycle et une semble
grimper sans limite, et l'heuristique explique pourquoi : un pas impair de l'application de
Collatz multiplie par 3 puis, en moyenne, divise par 4, pour un facteur net inférieur à 1,
tandis que le même décompte pour $ 5n+1 $ donne un facteur net supérieur à 1. Cet unique
accident arithmétique --- 3 est plus petit que 4, et 5 ne l'est pas --- est toute la raison
pour laquelle quiconque croit à la conjecture de Collatz, et toute démonstration devra
utiliser cet accident quelque part, exactement comme l'exige le test $ 3n-1 $ d'OPEN-18.
Notez aussi la symétrie de notre ignorance : on n'a jamais démontré d'aucune trajectoire
$ 5n+1 $ particulière qu'elle s'échappe vers l'infini, de même qu'il n'existe aucune
démonstration que toute trajectoire $ 3n+1 $ rentre au bercail. L'heuristique prédit avec
assurance dans les deux directions et ne démontre rien ni dans l'une ni dans l'autre.

\begin{refs}
\item Contexte : Conjecture de Syracuse (Collatz) --- Wikipédia (\url{https://fr.wikipedia.org/wiki/Conjecture_de_Syracuse})
\item Lecture : Jeffrey C. Lagarias, \og{} The 3x+1 problem: an annotated bibliography \fg{}, arXiv:math/0309224 (\url{https://arxiv.org/abs/math/0309224})
\item Lecture : Jeffrey C. Lagarias (dir.), \textit{The Ultimate Challenge: The 3x+1 Problem}, AMS (2010)
\end{refs}

\bigskip


\section*{La conjecture de Collatz}

\begin{problem}[{L'application $ 3n+1 $ : calculez, et voyez le problème --- Lycée}]
\textbf{OPEN-17.} \textbf{Problème.} Définissez l'application de Collatz sur les entiers strictement
positifs par
\[ T(n) = \begin{cases} n/2, & n \text{ pair},\\ 3n+1, & n \text{ impair}.\end{cases} \]
\begin{enumerate}
\item[(a)] Calculez la trajectoire complète de $ n = 27 $ jusqu'à ce qu'elle atteigne $ 1 $ ;
notez le nombre de pas qu'elle demande et la plus grande valeur atteinte. (Vous devriez
trouver que le voyage est bien plus long et bien plus haut que la taille de $ 27 $ ne le
laisse supposer --- c'est tout l'intérêt.)
\item[(b)] Définissez le \textit{temps d'arrêt total} de $ n $ comme le nombre de pas nécessaires
pour atteindre $ 1 $. Tabulez-le pour $ n = 1, \dots, 30 $ et trouvez quelle valeur
initiale inférieure à $ 30 $ met le plus longtemps.
\item[(c)] Expliquez avec vos propres mots pourquoi un calcul, aussi loin qu'il soit poussé, ne
pourra jamais démontrer la conjecture.
\end{enumerate}

\end{problem}

\noindent Lothar Collatz a consigné le problème dans ses carnets en 1937 ; il a
depuis été rattaché aux noms d'Ulam, Kakutani, Thwaites, Hasse et Syracuse, signe du
nombre de personnes qui l'ont redécouvert indépendamment. Kakutani plaisantait en disant
que c'était un complot pour ralentir la recherche mathématique américaine. La trajectoire
de $ 27 $ est la démonstration classique : elle grimpe jusqu'à $ 9232 $ et prend
$ 111 $ pas, sans raison visible. Toute valeur jusqu'à environ $ 2^{68} $, soit à peu
près $ 2{,}95 \times 10^{20} $, a été vérifiée comme atteignant $ 1 $ --- et cette
vérification ne nous dit rien sur le nombre suivant.

\begin{refs}
\item Contexte : Conjecture de Syracuse (Collatz) --- Wikipédia (\url{https://fr.wikipedia.org/wiki/Conjecture_de_Syracuse})
\item Lecture : Jeffrey C. Lagarias (dir.), \textit{The Ultimate Challenge: The 3x+1 Problem}, AMS (2010)
\item Lecture : Jeffrey C. Lagarias, \og{} The 3x+1 problem: An annotated bibliography \fg{}, arXiv:math/0309224 (\url{https://arxiv.org/abs/math/0309224})
\end{refs}

\bigskip

\begin{problem}[{L'application $ 3n-1 $ a des cycles --- trouvez-les --- Lycée}]
\textbf{OPEN-18.} \textbf{Problème.} Changez un signe. Définissez
\[ S(n) = \begin{cases} n/2, & n \text{ pair},\\ 3n-1, & n \text{ impair}.\end{cases} \]
\begin{enumerate}
\item[(a)] Montrez que $ S $ possède un cycle contenant $ 5 $, en calculant la trajectoire de
$ 5 $ jusqu'à son retour.
\item[(b)] Trouvez un second cycle, plus long, en partant de $ 17 $.
\item[(c)] Expliquez maintenant la portée de tout cela : la conjecture $ 3n+1 $ affirme
qu'aucun cycle autre que $ 1 \to 4 \to 2 \to 1 $ n'existe, et voici une application
presque identique où cette affirmation est \textbf{fausse}. Que vous apprend cela sur toute
démonstration proposée de Collatz qui fonctionnerait aussi bien pour $ 3n-1 $ ?
\end{enumerate}

\end{problem}

\noindent C'est l'exercice le plus utile de la page, et il ne coûte que de
l'arithmétique. Tout argument en faveur de la conjecture $ 3n+1 $ doit, quelque part,
utiliser une propriété qui distingue $ +1 $ de $ -1 $ ; s'il ne le fait pas, le même
argument démontrerait quelque chose de faux. Des générations de \og{} démonstrations \fg{}
amateurs de Collatz meurent exactement ici, et confronter votre propre raisonnement à
l'application $ 3n-1 $ est le contrôle d'honnêteté le moins coûteux des mathématiques.
Notez aussi que $ 3n-1 $ sur les entiers positifs revient au même que $ 3n+1 $ sur les
négatifs, raison pour laquelle le comportement de l'application de Collatz sur les entiers
négatifs fait partie du tableau standard.

\begin{refs}
\item Contexte : Conjecture de Syracuse (Collatz) --- Wikipédia (\url{https://fr.wikipedia.org/wiki/Conjecture_de_Syracuse})
\item Lecture : Lagarias, \textit{The Ultimate Challenge}, panorama introductif
\end{refs}

\bigskip

\begin{problem}[{Aucun cycle non trivial de longueur deux --- Licence}]
\textbf{OPEN-19.} \textbf{Problème.} Soit $ T $ l'application de Collatz d'OPEN-17.
\begin{enumerate}
\item[(a)] Démontrez que le seul cycle de l'application de Collatz dans lequel figure exactement
un nombre impair est $ 1 \to 4 \to 2 \to 1 $. Précisément : supposez $ n $ impair et
$ T^{k}(n) = n $, la trajectoire passant par exactement une valeur impaire ; montrez que
$ n = 1 $.
\item[(b)] Allez plus loin : supposez qu'un cycle contienne exactement deux nombres impairs
$ a $ et $ b $, de sorte que
\[ 3a+1 = 2^{i} b, \qquad 3b+1 = 2^{j} a \]
pour des entiers strictement positifs $ i, j $. Résolvez ce système en nombres entiers
et montrez que l'unique solution avec $ a, b \geq 1 $ est le cycle trivial.
\end{enumerate}

\end{problem}

\noindent C'est un morceau authentique, complet et digne d'un manuel, du problème
de Collatz, et il est entièrement à portée : un argument diophantien avec des majorations
élémentaires. Le faire vous montre à quoi ressemble un vrai résultat partiel sur Collatz,
et pourquoi le cas général est si difficile --- la même méthode étendue à $ k $ nombres
impairs produit une équation dont les solutions exigent des résultats profonds sur les
formes linéaires de logarithmes (le théorème de Baker), et c'est ainsi que sont
effectivement obtenues les bornes connues sur la longueur des cycles. Par de telles
méthodes, on sait que tout cycle non trivial serait immensément long --- les minorations
actuelles se comptent en centaines de millions de termes.

\begin{refs}
\item Contexte : Conjecture de Syracuse (Collatz) --- Wikipédia (\url{https://fr.wikipedia.org/wiki/Conjecture_de_Syracuse})
\item Contexte : Théorème de Baker --- Wikipédia (\url{https://fr.wikipedia.org/wiki/Th%C3%A9or%C3%A8me_de_Baker})
\item Lecture : Lagarias, \textit{The Ultimate Challenge}, chap. 1
\end{refs}

\bigskip

\begin{problem}[{Le théorème de densité de Terras --- Master}]
\textbf{OPEN-20.} \textbf{Problème.}
\begin{enumerate}
\item[(a)] Lisez et reconstituez le théorème de Riho Terras (1976) : presque tout entier
strictement positif $ n $, au sens de la densité naturelle, a un temps d'arrêt fini ---
c'est-à-dire qu'un certain itéré vérifie $ T^{k}(n) < n $.
\item[(b)] Identifiez l'heuristique probabiliste qui le sous-tend : une trajectoire de Collatz
\og{} aléatoire \fg{} multiplie par $ 3/2 $ aux pas impairs et par $ 1/2 $ aux pas pairs, de
sorte que le multiplicateur espéré par pas est $ \sqrt{3}/2 < 1 $.
\item[(c)] Énoncez précisément pourquoi cette heuristique n'est \textit{pas} une démonstration de
la conjecture.
\end{enumerate}

\end{problem}

\noindent Le résultat de Terras est le prototype de tout théorème sérieux sur
Collatz : il démontre que la conjecture vaut pour presque tous les entiers tout en
laissant ouverte la possibilité d'un ensemble clairsemé de contre-exemples --- et un seul
contre-exemple suffit. L'écart entre \og{} de densité un \fg{} et \og{} tous \fg{} est le problème tout
entier. L'heuristique est aussi la raison pour laquelle la plupart des mathématiciens
croient à la conjecture : en moyenne, les trajectoires devraient décroître. Croire et
démontrer sont deux activités différentes, et c'est avec Collatz que cette différence est
la plus visible.

\begin{refs}
\item Contexte : Conjecture de Syracuse, temps d'arrêt --- Wikipédia (\url{https://fr.wikipedia.org/wiki/Conjecture_de_Syracuse})
\item Article : R. Terras, \og{} A stopping time problem on the positive integers \fg{}, \textit{Acta Arithmetica} 30 (1976)
\item Lecture : Lagarias, \textit{The Ultimate Challenge}
\end{refs}

\bigskip

\begin{problem}[{Le théorème de Tao : presque toutes les orbites atteignent des valeurs presque bornées --- Recherche}]
\textbf{OPEN-21.} \textbf{Problème.} Lisez l'article de Terence Tao de 2019 et comprenez-en précisément
l'énoncé : pour toute fonction $ f $ tendant vers l'infini, si lentement soit-il,
presque tout $ n $ (au sens de la densité logarithmique) vérifie
$ \min_k T^{k}(n) < f(n) $. Répondez par écrit aux questions suivantes.
\begin{enumerate}
\item[(a)] Que signifie exactement \og{} presque tout \fg{} ici, et en quoi la densité logarithmique
diffère-t-elle de la densité naturelle ?
\item[(b)] Pourquoi cela reste-t-il en deçà de la conjecture complète --- que faudrait-il
renforcer ?
\item[(c)] Quel rôle la mesure invariante sur les entiers 3-adiques joue-t-elle dans
l'argument ?
\end{enumerate}

\end{problem}

\noindent C'est le résultat général le plus fort connu sur Collatz, et il vient
de l'un des meilleurs analystes vivants ; Tao a décrit le problème comme \og{} notoirement
difficile \fg{} et largement rétif aux techniques disponibles. La méthode de l'article ---
traiter l'application de Collatz de façon probabiliste et trouver une mesure invariante ---
est la frontière actuelle, et le lire attentivement constitue une véritable formation de
niveau recherche sur la manière dont les techniques analytiques modernes se braquent sur
une question d'apparence élémentaire. Notez l'expression \og{} presque bornées \fg{} : même ce
résultat héroïque ne peut exclure qu'une orbite rare s'échappe vers l'infini.

\begin{refs}
\item Article : Terence Tao, \og{} Almost all orbits of the Collatz map attain almost bounded values \fg{} (2019), arXiv:1909.03562 (\url{https://arxiv.org/abs/1909.03562})
\item Lecture : le billet de blog de Tao expliquant le résultat (en anglais) (\url{https://terrytao.wordpress.com/2019/09/10/almost-all-collatz-orbits-attain-almost-bounded-values/})
\item Contexte : Conjecture de Syracuse (Collatz) --- Wikipédia (\url{https://fr.wikipedia.org/wiki/Conjecture_de_Syracuse})
\end{refs}

\bigskip

\begin{problem}[{Le FRACTRAN de Conway et l'indécidabilité --- Master}]
\textbf{OPEN-22.} \textbf{Problème.} En FRACTRAN, le langage de Conway, un programme est une liste de
fractions et, à partir d'un état $ n $, on multiplie par la première fraction de la
liste qui donne un entier.
\begin{enumerate}
\item[(a)] Apprenez FRACTRAN, et démontrez qu'il est Turing-complet.
\item[(b)] Comprenez le théorème de Conway de 1972 : la généralisation naturelle du problème de
Collatz --- les applications définies par morceaux par $ n \mapsto a_i n + b_i $ selon
$ n \bmod d $ --- est \textbf{indécidable} ; il n'existe aucun algorithme qui décide, pour
une telle application et une valeur initiale arbitraires, si l'orbite atteint $ 1 $.
\item[(c)] Expliquez soigneusement ce que cela dit, et ne dit pas, de la conjecture $ 3n+1 $
originale, qui est une instance fixée unique et n'est donc pas couverte par le résultat
d'indécidabilité.
\end{enumerate}

\end{problem}

\noindent Le résultat de John Horton Conway est le fait structurel le plus profond
connu sur cette famille de problèmes, et il reformule la remarque d'Erdős : si Collatz
résiste, c'est peut-être qu'il se situe juste en dessous d'un seuil de calcul universel.
Une application de Collatz généralisée peut simuler n'importe quel programme
informatique ; s'interroger sur son comportement à long terme, c'est donc poser le
problème de l'arrêt. La question (c) est le point intellectuellement délicat que bien des
récits de vulgarisation manquent --- l'indécidabilité d'une famille ne rend aucun membre
particulier indémontrable, et la conjecture $ 3n+1 $ pourrait encore avoir une
démonstration de deux pages que personne n'a trouvée.

\begin{refs}
\item Contexte : FRACTRAN --- Wikipédia (\url{https://fr.wikipedia.org/wiki/FRACTRAN})
\item Contexte : Généralisations indécidables --- Wikipédia (\url{https://fr.wikipedia.org/wiki/Conjecture_de_Syracuse})
\item Article : J. H. Conway, \og{} Unpredictable iterations \fg{} (1972), réédité dans Lagarias, \textit{The Ultimate Challenge}
\item Voir aussi : Mathématiques discrètes et informatique (\url{applied-discrete-cs.html})
\end{refs}

\bigskip

\begin{problem}[{Le temps d'arrêt total de 97 --- Lycée, short}]
\textbf{OPEN-47.} \textbf{Problème.} À l'aide de l'application de Collatz $ T $
d'OPEN-17, calculez le temps d'arrêt total de $ n = 97 $ --- le nombre de pas nécessaires
pour atteindre $ 1 $ --- et notez la plus grande valeur que sa trajectoire atteint en
chemin. Comparez ces deux nombres à ceux qui correspondent à $ n = 27 $, et dites ce que
la comparaison suggère quant à la prédiction d'une trajectoire à partir de la taille de sa
valeur initiale.
\end{problem}

\noindent Le faire une fois, à la main ou en cinq lignes de code, vaut mieux que
lire n'importe quelle description du problème. Deux choses deviennent visibles : à quelle
hauteur une trajectoire peut errer au-dessus de son point de départ, et comment des
valeurs initiales différentes sont canalisées vers des chemins communs, de sorte que des
nombres tout à fait étrangers l'un à l'autre finissent par descendre la même échelle
jusqu'à $ 1 $. La suite des temps d'arrêt totaux est cataloguée sous OEIS A006577, et
son graphe --- hérissé, sans structure, n'obéissant à aucune formule que quiconque ait
trouvée --- est la raison pour laquelle Erdős disait que les mathématiques ne sont
peut-être pas prêtes pour ce problème.

\begin{refs}
\item Contexte : Conjecture de Syracuse (Collatz) --- Wikipédia (\url{https://fr.wikipedia.org/wiki/Conjecture_de_Syracuse})
\item Données : OEIS A006577 (\url{https://oeis.org/A006577}) --- nombre de pas pour atteindre 1 dans le problème $ 3x+1 $
\end{refs}

\bigskip


\section*{Problèmes ouverts classiques}

\begin{problem}[{La conjecture de Goldbach --- Recherche}]
\textbf{OPEN-23.} \textbf{Problème.} \textbf{Démontrez que tout entier pair supérieur à
$ 2 $ est une somme de deux nombres premiers.} Statut : ouvert. La conjecture de
Goldbach \textit{faible} (ternaire) --- tout entier impair supérieur à $ 5 $ est une somme
de trois nombres premiers --- a été démontrée par Harald Helfgott en 2013 ; sa démonstration
est acceptée, même si la publication formelle du manuscrit complet en revue a longtemps
tardé. Chen Jingrun a démontré en 1973 que tout entier pair assez grand est la somme d'un
nombre premier et d'un nombre ayant au plus deux facteurs premiers.
\end{problem}

\noindent Christian Goldbach a écrit la conjecture à Leonhard Euler dans une
lettre du 7 juin 1742 ; Euler a répondu qu'il la tenait pour certaine, sans pouvoir la
démontrer. Elle a été vérifiée au-delà de $ 4 \times 10^{18} $. L'écart entre la forme
faible et la forme forte est instructif : trois nombres premiers laissent assez de marge
pour que la méthode du cercle réussisse, deux n'en laissent pas assez, et aucune technique
connue ne comble cette différence. C'est le motif récurrent de la théorie additive des
nombres --- les problèmes deviennent traitables exactement quand on vous accorde un degré de
liberté de plus que vous n'en vouliez.

\begin{refs}
\item Contexte : Conjecture de Goldbach --- Wikipédia (\url{https://fr.wikipedia.org/wiki/Conjecture_de_Goldbach})
\item Article : H. A. Helfgott, \og{} The ternary Goldbach conjecture is true \fg{} (2013), arXiv:1312.7748 (\url{https://arxiv.org/abs/1312.7748})
\item Voir aussi : Théorie des nombres (\url{pure-number-theory.html})
\end{refs}

\bigskip

\begin{problem}[{Les nombres premiers jumeaux, et la révolution des écarts bornés --- Recherche}]
\textbf{OPEN-24.} \textbf{Problème.} La question voisine des écarts \textit{bornés} a été percée : Yitang
Zhang a démontré en 2013 que
\[ \liminf_{n \to \infty}(p_{n+1} - p_n) < 7 \times 10^{7}, \]
et après les raffinements de James Maynard, Terence Tao et la collaboration Polymath8, la
borne s'établit aujourd'hui à $ 246 $.
\begin{enumerate}
\item[(a)] \textbf{Démontrez qu'il existe une infinité de nombres premiers $ p $ tels que
$ p + 2 $ soit également premier.} Statut : ouvert.
\item[(b)] Comprenez la différence entre \og{} un écart au plus égal à 246 une infinité de fois \fg{} et
\og{} un écart exactement égal à 2 une infinité de fois \fg{}, et pourquoi les méthodes de crible
employées ne peuvent pas atteindre $ 2 $.
\end{enumerate}

\end{problem}

\noindent L'article de Zhang est l'une des grandes histoires des mathématiques
modernes : c'était un chargé de cours peu connu, de près de cinquante-huit ans, qui avait
travaillé des années en marge du courant dominant de la recherche, et il a soumis une
démonstration complète d'un résultat que les spécialistes croyaient éloigné de plusieurs
décennies. La borne de soixante-dix millions était, comme chacun l'a immédiatement
compris, sans importance --- l'essentiel était qu'une borne finie existât, quelle qu'elle
fût. En quelques mois, un projet Polymath en ligne, massivement collaboratif, l'avait
ramenée à quelques centaines. La barrière à $ 246 $, ou à $ 6 $ sous la conjecture non
démontrée d'Elliott--Halberstam, est une limitation réelle de la théorie du crible connue
sous le nom de problème de parité, et parvenir à $ 2 $ demandera une idée que personne
n'a.

\begin{refs}
\item Contexte : Nombres premiers jumeaux --- Wikipédia (\url{https://fr.wikipedia.org/wiki/Nombres_premiers_jumeaux})
\item Article : Yitang Zhang, \og{} Bounded gaps between primes \fg{}, \textit{Annals of Mathematics} 179 (2014)
\item Article : James Maynard, \og{} Small gaps between primes \fg{}, \textit{Annals of Mathematics} 181 (2015), arXiv:1311.4600 (\url{https://arxiv.org/abs/1311.4600})
\item Contexte : Projet Polymath --- Wikipédia (\url{https://fr.wikipedia.org/wiki/Projet_Polymath})
\end{refs}

\bigskip

\begin{problem}[{Les quatre problèmes de Landau --- Recherche}]
\textbf{OPEN-25.} \textbf{Problème.} Au Congrès international de 1912, Edmund Landau a énuméré quatre
problèmes sur les nombres premiers qu'il jugeait \og{} inattaquables dans l'état actuel de la
science \fg{}. \textbf{Les quatre restent ouverts en 2026.} Pour chacun, notez le meilleur
résultat partiel connu et énoncez exactement ce qui lui manque.
\begin{enumerate}
\item[(a)] Goldbach : tout $ n > 2 $ pair est une somme de deux nombres premiers.
\item[(b)] Nombres premiers jumeaux : il existe une infinité de $ p $ tels que $ p+2 $ soit
premier.
\item[(c)] Conjecture de Legendre : il y a toujours un nombre premier entre $ n^2 $ et
$ (n+1)^2 $.
\item[(d)] Existe-t-il une infinité de nombres premiers de la forme $ n^2 + 1 $ ?
\end{enumerate}

\end{problem}

\noindent Un siècle et une décennie plus tard, le jugement de Landau tient
toujours, ce qui en dit long sur la difficulté des questions additives concernant les
nombres premiers. Du terrain a été gagné : le postulat de Bertrand donne un nombre premier
entre $ n $ et $ 2n $ (exercice facile), et Baker--Harman--Pintz ont démontré qu'il y a
un nombre premier dans $ [x, x + x^{0.525}] $ pour $ x $ grand --- mais un écart de
$ x^{0.525} $ est bien plus grand que les quelque $ 2\sqrt{n} $ qu'exige la conjecture
de Legendre, et même l'hypothèse de Riemann ne le refermerait pas. Iwaniec a démontré
qu'une infinité de $ n^2+1 $ ont au plus deux facteurs premiers. Rassembler en un seul
endroit ces quatre meilleurs résultats connus est une véritable formation à la forme du
sujet.

\begin{refs}
\item Contexte : Problèmes de Landau --- Wikipédia (\url{https://fr.wikipedia.org/wiki/Probl%C3%A8mes_de_Landau})
\item Contexte : Conjecture de Legendre --- Wikipédia (\url{https://fr.wikipedia.org/wiki/Conjecture_de_Legendre})
\item Lecture : G. H. Hardy \& E. M. Wright, \textit{An Introduction to the Theory of Numbers}
\end{refs}

\bigskip

\begin{problem}[{Les nombres parfaits impairs --- Recherche}]
\textbf{OPEN-26.} \textbf{Problème.} Un nombre est \textit{parfait} s'il est égal à la somme de ses diviseurs
propres, autrement dit si $ \sigma(n) = 2n $.
\begin{enumerate}
\item[(a)] Démontrez d'abord la moitié classique facile : tout nombre parfait \textit{pair} a la
forme d'Euclide--Euler $ 2^{p-1}(2^p - 1) $ avec $ 2^p - 1 $ premier.
\item[(b)] \textbf{Existe-t-il un nombre parfait impair ?} Statut : ouvert. D'énormes contraintes
sont connues --- tout nombre parfait impair devrait dépasser $ 10^{1500} $, avoir au moins
dix facteurs premiers distincts et satisfaire une longue liste de conditions de congruence
--- mais aucune démonstration de non-existence.
\end{enumerate}

\end{problem}

\noindent C'est peut-être le plus ancien problème ouvert des mathématiques : les
Grecs connaissaient $ 6 $ et $ 28 $, Euclide a démontré la suffisance de sa formule
vers 300 av. J.-C., et Euler en a démontré la nécessité pour les nombres pairs deux mille
ans plus tard. Le cas impair résiste depuis. C'est un bel exemple de problème où une
montagne de conditions nécessaires s'est accumulée sans que personne trouve la
contradiction --- chaque contrainte publiée réduit la recherche et aucune ne la referme.
Notez la possibilité honnête, si peu ragoûtante soit-elle : qu'un nombre parfait impair
existe et soit simplement trop grand pour qu'on le trouve.

\begin{refs}
\item Contexte : Nombre parfait --- Wikipédia (\url{https://fr.wikipedia.org/wiki/Nombre_parfait})
\item Contexte : Théorème d'Euclide-Euler --- Wikipédia (\url{https://fr.wikipedia.org/wiki/Th%C3%A9or%C3%A8me_d%27Euclide-Euler})
\item Voir aussi : Théorie des nombres (\url{pure-number-theory.html})
\end{refs}

\bigskip

\begin{problem}[{La conjecture abc, et une controverse --- Recherche}]
\textbf{OPEN-27.} \textbf{Problème.} Pour des entiers strictement positifs premiers entre
eux vérifiant $ a + b = c $, notons $ \operatorname{rad}(n) $ le produit des nombres
premiers distincts divisant $ n $. La conjecture de Masser et Oesterlé affirme que, pour
tout $ \varepsilon > 0 $, il n'existe qu'un nombre fini de tels triplets avec
\[ c > \operatorname{rad}(abc)^{1+\varepsilon} . \]
Statut : \textbf{contesté}. Shinichi Mochizuki a publié en 2021, dans \textit{PRIMS}, une
démonstration revendiquée passant par la théorie de Teichmüller inter-universelle, mais
cette démonstration n'est pas acceptée par la communauté mathématique au sens large :
Peter Scholze et Jakob Stix ont publié en 2018 une objection détaillée identifiant une
étape précise (dans la démonstration du corollaire 3.12) qu'ils jugent irrémédiablement
défectueuse, et Mochizuki conteste leur objection. En 2026, il n'existe aucun consensus
selon lequel abc serait démontrée.
\end{problem}

\noindent La conjecture est d'une puissance saisissante --- elle implique le dernier
théorème de Fermat pour les grands exposants en quelques lignes, ainsi qu'une longue liste
d'autres résultats --- et c'est précisément pour cela que son statut importe tant et que la
controverse a été si douloureuse. Cette entrée figure ici en partie comme exercice
d'honnêteté intellectuelle : les mathématiques sont une activité humaine, la vérification
est un processus social autant que logique, et \og{} une démonstration a été publiée \fg{} n'est
pas toujours la même chose que \og{} le théorème est acquis \fg{}. Lisez les deux camps avant de
vous forger une opinion, et remarquez avec quel soin chacun expose sa position.

\begin{refs}
\item Contexte : Conjecture abc --- Wikipédia (\url{https://fr.wikipedia.org/wiki/Conjecture_abc})
\item Contexte : Théorie de Teichmüller inter-universelle --- Wikipédia (en anglais) (\url{https://en.wikipedia.org/wiki/Inter-universal_Teichm%C3%BCller_theory})
\item Lecture : Peter Scholze \& Jakob Stix, \og{} Why abc is still a conjecture \fg{} (2018)
\end{refs}

\bigskip

\begin{problem}[{L'irrationalité de $ \gamma $, et celle de $ \zeta(5) $ --- Recherche}]
\textbf{OPEN-28.} \textbf{Problème.} La constante d'Euler--Mascheroni est
\[ \gamma = \lim_{n \to \infty}\left( \sum_{k=1}^{n} \frac{1}{k} - \log n \right) \approx 0.5772156649 . \]
\begin{enumerate}
\item[(a)] \textbf{Démontrez que $ \gamma $ est irrationnel.} Statut : ouvert --- on ne sait même
pas si $ \gamma $ est irrationnel, encore moins s'il est transcendant.
\item[(b)] Apéry a démontré en 1978 que $ \zeta(3) $ est irrationnel. \textbf{Démontrez que
$ \zeta(5) $ est irrationnel.} Statut : ouvert ; on sait qu'au moins l'un des nombres
$ \zeta(5), \zeta(7), \zeta(9), \zeta(11) $ est irrationnel, mais pas lequel.
\end{enumerate}

\end{problem}

\noindent $ \gamma $ a été calculé avec plus de mille milliards de décimales et
apparaît partout en analyse et en théorie des nombres, et nous ne savons pas démontrer à
son sujet le fait qualitatif le plus élémentaire. La démonstration d'Apéry pour
$ \zeta(3) $, en 1978, est légendaire par son accueil : il a présenté une suite
d'identités d'apparence non motivée à un auditoire sceptique, dont plusieurs membres sont
partis, et il a fallu des semaines de vérification avant qu'on admette qu'il s'agissait
d'une démonstration réelle et élémentaire de quelque chose que tout le monde croyait
éloigné de plusieurs décennies. Cette histoire vaut d'être connue quand votre propre
argument a l'air bizarre.

\begin{refs}
\item Contexte : Constante d'Euler-Mascheroni --- Wikipédia (\url{https://fr.wikipedia.org/wiki/Constante_d%27Euler-Mascheroni})
\item Contexte : Théorème d'Apéry --- Wikipédia (\url{https://fr.wikipedia.org/wiki/Th%C3%A9or%C3%A8me_d%27Ap%C3%A9ry})
\item Lecture : Alfred van der Poorten, \og{} A proof that Euler missed\dots{} Apéry's proof of the irrationality of $\zeta$(3) \fg{} (1979)
\item Voir aussi : Analyse (\url{applied-calculus.html}), L'art de la démonstration (\url{proofs.html})
\end{refs}

\bigskip

\begin{problem}[{Le problème du carré inscrit --- Recherche}]
\textbf{OPEN-29.} \textbf{Problème.} \textbf{Démontrez que toute courbe de Jordan --- toute
courbe fermée simple continue du plan --- contient quatre points formant les sommets d'un
carré.} Statut : ouvert dans le cas général ; posé par Otto Toeplitz en 1911. Le
résultat est connu pour les courbes convexes, pour les courbes analytiques par morceaux,
et (Greene--Lobb, 2020) pour les courbes lisses dans le cas des \textit{rectangles} inscrits
de tout rapport d'aspect. L'obstruction tient entièrement aux courbes qui ne sont que
continues, et donc possiblement sauvages.
\end{problem}

\noindent Tracez n'importe quelle boucle fermée qui ne se recoupe pas, aussi
déchiquetée soit-elle. Pouvez-vous toujours y trouver quatre points formant un carré
parfait ? Tout le monde croit que oui et personne ne sait le démontrer, ce qui en fait le
problème ouvert le plus accessible de la géométrie --- l'énoncé ne demande aucun vocabulaire
au-delà du lycée, et la difficulté est réelle et profonde. La démonstration de Joshua
Greene et Andrew Lobb en 2020 pour les courbes lisses a utilisé la géométrie symplectique,
un outil de la physique mathématique, pour régler la version rectangulaire ; c'est dans le
saut du lisse au continu que vivent les monstres fractals.

\begin{refs}
\item Contexte : Problème du carré inscrit --- Wikipédia (\url{https://fr.wikipedia.org/wiki/Probl%C3%A8me_du_carr%C3%A9_inscrit})
\item Article : Joshua Greene \& Andrew Lobb, \og{} The rectangular peg problem \fg{} (2020), arXiv:2005.09193 (\url{https://arxiv.org/abs/2005.09193})
\item Voir aussi : Géométrie (\url{pure-geometry.html}), Topologie (\url{pure-topology.html})
\end{refs}

\bigskip

\begin{problem}[{Le nombre chromatique du plan --- Recherche}]
\textbf{OPEN-30.} \textbf{Problème.} Coloriez chaque point du plan de sorte que deux points à distance
exactement $ 1 $ n'aient jamais la même couleur, et soit $ \chi $ le nombre minimal de
couleurs nécessaires. Statut : ouvert, mais resserré. Pendant soixante ans on savait
seulement que la réponse était dans $ \{4,5,6,7\} $ ; en 2018 Aubrey de Grey a trouvé un
graphe de distance unité à $ 1581 $ sommets qui ne peut pas être colorié avec 4
couleurs, ce qui a relevé la borne inférieure, d'où $ 5 \leq \chi \leq 7 $.
\begin{enumerate}
\item[(a)] Démontrez vous-même les deux bornes classiques faciles : le fuseau de Moser force
$ \chi \geq 4 $, et un pavage hexagonal donne $ \chi \leq 7 $.
\item[(b)] \textbf{Déterminez $ \chi $.}
\end{enumerate}

\end{problem}

\noindent Le problème de Hadwiger--Nelson a été posé vers 1950 et il est célèbre
pour le peu qui sépare l'amateur du professionnel : la borne supérieure est un dessin
d'hexagones que n'importe qui peut tracer, la borne inférieure un petit graphe que
n'importe qui peut vérifier, et la vérité se trouve quelque part dans un intervalle que
personne n'a refermé en soixante-dix ans. De Grey --- plus connu pour ses recherches sur le
vieillissement, et travaillant là-dessus en électron libre --- a annoncé son amélioration
dans un article qu'un projet Polymath a ensuite vérifié et simplifié. C'est la
démonstration récente la plus claire que cette page n'est pas réservée aux
professionnels.

\begin{refs}
\item Contexte : Problème de Hadwiger-Nelson --- Wikipédia (\url{https://fr.wikipedia.org/wiki/Probl%C3%A8me_de_Hadwiger-Nelson})
\item Article : Aubrey D. N. J. de Grey, \og{} The chromatic number of the plane is at least 5 \fg{} (2018), arXiv:1804.02385 (\url{https://arxiv.org/abs/1804.02385})
\item Voir aussi : Combinatoire et théorie des graphes (\url{pure-combinatorics.html})
\end{refs}

\bigskip

\begin{problem}[{La conjecture des familles stables par unions --- Recherche}]
\textbf{OPEN-31.} \textbf{Problème.} Une famille finie d'ensembles est \textit{stable par
union} si la réunion de deux de ses membres en est encore un membre. \textbf{Démontrez que
toute famille stable par union (autre que $ \{\emptyset\} $) contient un élément
appartenant à au moins la moitié de ses ensembles.} Statut : ouvert. En 2022, Justin
Gilmer a donné la première démonstration d'une minoration par une constante, à l'aide d'un
argument d'entropie, montrant qu'un certain élément se trouve dans au moins une fraction
$ 0.01 $ des ensembles ; sa méthode a été rapidement affinée par plusieurs équipes
jusqu'à $ (3-\sqrt5)/2 \approx 0.38 $. La constante conjecturée $ 1/2 $ n'est
toujours pas démontrée.
\end{problem}

\noindent Péter Frankl a posé cette question en 1979 et, pendant quarante ans,
personne n'a su démontrer \textit{la moindre} fraction constante, ce qui est stupéfiant pour
un énoncé aussi simple. La percée de Gilmer est un bel exemple de transfert de technique :
il a appliqué la théorie de l'information, en mesurant l'entropie de choix aléatoires
d'ensembles, à un problème d'apparence purement combinatoire, et la barrière est tombée
quelques semaines après la parution de sa prépublication. On sait aujourd'hui que l'écart
restant entre $ 0.38 $ et $ 0.5 $ est hors de portée de cette méthode particulière.

\begin{refs}
\item Contexte : Conjecture des familles stables par unions --- Wikipédia (\url{https://fr.wikipedia.org/wiki/Conjecture_des_familles_stables_par_unions})
\item Article : Justin Gilmer, \og{} A constant lower bound for the union-closed sets conjecture \fg{} (2022), arXiv:2211.09055 (\url{https://arxiv.org/abs/2211.09055})
\item Voir aussi : Combinatoire et théorie des graphes (\url{pure-combinatorics.html})
\end{refs}

\bigskip

\begin{problem}[{Petits nombres de Ramsey : que vaut $ R(5,5) $ ? --- Recherche}]
\textbf{OPEN-32.} \textbf{Problème.} $ R(5,5) $ est le plus petit $ n $ tel que tout coloriage en rouge
et bleu des arêtes de $ K_n $ contienne un $ K_5 $ monochrome. Statut : ouvert. On
sait que $ 43 \leq R(5,5) \leq 46 $, la borne supérieure étant due à Angeltveit et
McKay (2024).
\begin{enumerate}
\item[(a)] Démontrez vous-même le petit cas classique $ R(3,3) = 6 $ --- à la fois le coloriage
montrant $ R(3,3) > 5 $ et l'argument des tiroirs montrant $ R(3,3) \leq 6 $.
\item[(b)] \textbf{Déterminez $ R(5,5) $.}
\end{enumerate}

\end{problem}

\noindent La célèbre remarque d'Erdős dit les choses mieux que n'importe quel
commentaire : si des extraterrestres exigeaient la valeur de $ R(5,5) $ sous peine de
détruire la Terre, nous devrions mobiliser tous nos ordinateurs et tous nos mathématiciens
pour tenter de la trouver --- mais s'ils demandaient $ R(6,6) $, nous ferions mieux
d'essayer de détruire les extraterrestres. L'espace de recherche est astronomiquement
grand et la force brute est sans espoir ; les bornes ci-dessus proviennent d'un
raisonnement combinatoire profond assorti d'énormes calculs. Du côté des
\textit{asymptotiques}, en revanche, il y a eu du mouvement récemment : Campos, Griffiths,
Morris et Sahasrabudhe ont obtenu en 2023 la première amélioration exponentielle de la
borne supérieure de Ramsey diagonale, après quatre-vingt-huit ans.

\begin{refs}
\item Contexte : Théorème de Ramsey --- Wikipédia (\url{https://fr.wikipedia.org/wiki/Th%C3%A9or%C3%A8me_de_Ramsey})
\item Lecture : Stanisław Radziszowski, \og{} Small Ramsey Numbers \fg{}, panorama dynamique de l'\textit{Electronic Journal of Combinatorics} --- la table de référence des bornes actuelles
\item Article : Campos, Griffiths, Morris \& Sahasrabudhe, \og{} An exponential improvement for diagonal Ramsey \fg{} (2023), arXiv:2303.09521 (\url{https://arxiv.org/abs/2303.09521})
\end{refs}

\bigskip

\begin{problem}[{La conjecture du coureur solitaire --- Recherche}]
\textbf{OPEN-33.} \textbf{Problème.} Supposons que $ n $ coureurs partent ensemble sur
une piste circulaire de circonférence $ 1 $ et courent indéfiniment à des vitesses
constantes deux à deux distinctes. Un coureur est \textit{solitaire} à un instant où tous
les autres sont à une distance d'au moins $ 1/n $ le long de la piste. \textbf{Démontrez que
chaque coureur est solitaire à un certain moment.} Statut : ouvert dans le cas général ;
démontré pour $ n \leq 7 $ coureurs.
\end{problem}

\noindent La conjecture a été formulée indépendamment par Jörg Wills (1967) et
Thomas Cusick (1973) dans le langage de l'approximation diophantienne, et elle doit son nom
charmant à Luis Goddyn, en 1998. Chaque nouvelle valeur de $ n $ a exigé un travail
substantiellement neuf, signature d'un problème pour lequel aucune idée unificatrice n'a
encore été trouvée. C'est un bon exemple de la manière dont une question sur l'approximation
simultanée de plusieurs nombres irrationnels peut être habillée sous une forme qu'un enfant
peut se représenter --- et du peu que cela aide.

\begin{refs}
\item Contexte : Conjecture du coureur solitaire --- Wikipédia (\url{https://fr.wikipedia.org/wiki/Conjecture_du_coureur_solitaire})
\item Contexte : Approximation diophantienne --- Wikipédia (\url{https://fr.wikipedia.org/wiki/Approximation_diophantienne})
\end{refs}

\bigskip

\begin{problem}[{Le problème du sous-espace invariant --- Recherche}]
\textbf{OPEN-34.} \textbf{Problème.} \textbf{Tout opérateur linéaire borné sur un espace de
Hilbert complexe séparable de dimension infinie possède-t-il un sous-espace invariant
fermé non trivial ?} Statut : ouvert pour les espaces de Hilbert. Pour les espaces de
Banach généraux la réponse est \textit{non} --- Per Enflo a construit un contre-exemple
(annoncé en 1975, publié en 1987) et Charles Read en a produit d'autres. Notez que des
résolutions revendiquées du cas hilbertien ont été annoncées ces dernières années, y
compris par Enflo lui-même en 2023 ; elles n'ont pas été acceptées par la communauté, et
en 2026 le problème doit être considéré comme ouvert.
\end{problem}

\noindent C'est la question ouverte centrale de la théorie des opérateurs, et son
énoncé est trompeusement simple : tout opérateur préserve-t-il quelque chose ? En dimension
finie, les vecteurs propres règlent la question immédiatement. Le cas hilbertien de
dimension infinie a absorbé des efforts considérables, et l'on sait beaucoup de choses pour
des classes particulières --- opérateurs compacts (Aronszajn--Smith), opérateurs normaux,
opérateurs polynomialement compacts (Bernstein--Robinson, par l'analyse non standard). Cette
entrée sert aussi de leçon de lecture attentive des annonces : une prépublication n'est pas
un théorème tant que la communauté ne l'a pas vérifiée.

\begin{refs}
\item Contexte : Problème du sous-espace invariant --- Wikipédia (en anglais) (\url{https://en.wikipedia.org/wiki/Invariant_subspace_problem})
\item Lecture : John B. Conway, \textit{A Course in Functional Analysis}
\item Voir aussi : Analyse fonctionnelle (\url{pure-functional-analysis.html})
\end{refs}

\bigskip

\begin{problem}[{La conjecture jacobienne --- Recherche}]
\textbf{OPEN-35.} \textbf{Problème.} Soit $ F : \mathbb{C}^n \to \mathbb{C}^n $ une
application polynomiale dont le déterminant jacobien est une constante non nulle.
\textbf{Démontrez que $ F $ est inversible, d'inverse polynomial.} Statut : ouvert pour
tout $ n \geq 2 $. Des réductions connues montrent qu'il suffit de traiter les
applications de la forme particulière $ F(x) = x + H(x) $ avec $ H $ cubique homogène
--- et même ce cas n'est pas résolu.
\end{problem}

\noindent Ott-Heinrich Keller a posé cette question en 1939 et elle est devenue
tristement célèbre par le nombre de fausses démonstrations publiées dans les deux sens ;
Shreeram Abhyankar l'appelait \og{} le cimetière des géomètres algébristes \fg{}. L'énoncé se situe
à un curieux carrefour d'algèbre, de géométrie et d'analyse : l'hypothèse est une condition
d'inversibilité purement locale et la conclusion est globale et algébrique, et toute la
difficulté tient à ce que les polynômes sur $ \mathbb{C} $ refusent d'obéir à l'intuition
que le théorème d'inversion locale donne sur $ \mathbb{R} $. Notez que la conjecture est
\textit{fausse} en caractéristique positive, avertissement qu'il vaut la peine de retenir.

\begin{refs}
\item Contexte : Conjecture jacobienne --- Wikipédia (\url{https://fr.wikipedia.org/wiki/Conjecture_jacobienne})
\item Lecture : Arno van den Essen, \textit{Polynomial Automorphisms and the Jacobian Conjecture}
\end{refs}

\bigskip

\begin{problem}[{La conjecture d'Erdős--Straus --- Recherche}]
\textbf{OPEN-36.} \textbf{Problème.} Erdős et Straus ont conjecturé que, pour tout entier $ n \geq 2 $,
il existe des entiers strictement positifs $ x, y, z $ tels que
\[ \frac{4}{n} = \frac{1}{x} + \frac{1}{y} + \frac{1}{z} . \]
Statut : ouvert. Vérifiée par ordinateur pour tout $ n $ jusqu'à $ 10^{17} $ et
au-delà, et connue pour tout $ n $ hors de certaines classes résiduelles modulo
$ 840 $ --- les cas difficiles sont essentiellement les $ n $ premiers congrus à
$ 1, 11^2, 13^2, 17^2, 19^2 $ ou $ 23^2 $ modulo $ 840 $.
\begin{enumerate}
\item[(a)] Échauffez-vous en trouvant à la main des représentations pour $ n = 5, 7, 9 $.
\item[(b)] \textbf{Démontrez la conjecture pour tout $ n \geq 2 $.}
\end{enumerate}

\end{problem}

\noindent Erdős et Ernst Straus ont posé cette question vers 1948. Les fractions
égyptiennes --- sommes de fractions unitaires distinctes --- comptent parmi les objets les
plus anciens des mathématiques écrites, puisqu'elles apparaissent dans le papyrus Rhind
vers 1650 av. J.-C. ; c'est donc une question moderne dans une langue antique. Ce qui la
rend typique d'Erdős, c'est qu'elle est élémentaire à énoncer, trivialement vérifiable sur
n'importe quelle instance, et complètement rétive dans le cas général : les résultats
connus viennent tous de systèmes de congruences couvrants, et aucun système couvrant ne
traite les classes restantes.

\begin{refs}
\item Contexte : Conjecture d'Erdős-Straus --- Wikipédia (\url{https://fr.wikipedia.org/wiki/Conjecture_d%27Erd%C5%91s-Straus})
\item Contexte : Fraction égyptienne --- Wikipédia (\url{https://fr.wikipedia.org/wiki/Fraction_%C3%A9gyptienne})
\item Lecture : Erdős Problems (\url{https://www.erdosproblems.com/}) --- le catalogue tenu par Thomas Bloom des problèmes d'Erdős et de leur statut
\end{refs}

\bigskip

\begin{problem}[{Goldbach à la main, de 4 à 60 --- Lycée, short}]
\textbf{OPEN-48.} \textbf{Problème.} Pour tout nombre pair $ n $ avec
$ 4 \leq n \leq 60 $, écrivez $ n $ comme somme de deux nombres premiers et notez
combien de représentations non ordonnées $ n = p + q $ avec $ p \leq q $ vous trouvez.
Décrivez ensuite ce qu'il advient de ce décompte quand $ n $ grandit, et énoncez
clairement pourquoi aucune quantité de vérifications de ce genre ne règle OPEN-23.
\end{problem}

\noindent C'est le calcul que faisaient Goldbach et Euler en 1742, et il demande
une vingtaine de minutes avec une liste des nombres premiers jusqu'à $ 60 $. Le nombre
de représentations ne se contente pas de rester positif, il croît --- son graphique est
connu sous le nom de comète de Goldbach, et sa forme est prédite assez précisément par
l'heuristique de la méthode du cercle de Hardy--Littlewood. Cette croissance explique
pourquoi presque tout le monde croit à la conjecture, et aussi pourquoi personne ne sait
la démontrer : un décompte asymptotique grand en moyenne ne dit rien sur la possibilité
qu'un nombre pair isolé soit malchanceux.

\begin{refs}
\item Contexte : Conjecture de Goldbach --- Wikipédia (\url{https://fr.wikipedia.org/wiki/Conjecture_de_Goldbach})
\item Contexte : Comète de Goldbach --- Wikipédia (\url{https://fr.wikipedia.org/wiki/Com%C3%A8te_de_Goldbach})
\end{refs}

\bigskip

\begin{problem}[{Radicaux de triplets abc --- Licence, short}]
\textbf{OPEN-49.} \textbf{Problème.} Pour chaque triplet d'entiers premiers entre eux
avec $ a + b = c $ de la liste
\[ (1,8,9), \quad (5,27,32), \quad (1,80,81), \quad (3,125,128) \]
calculez $ \operatorname{rad}(abc) $, le produit des nombres premiers distincts divisant
$ abc $, et déterminez lesquels des quatre vérifient
$ c > \operatorname{rad}(abc) $. Calculez ensuite la \textit{qualité}
$ q = \log c / \log \operatorname{rad}(abc) $ de chaque triplet et énoncez ce que la
conjecture abc d'OPEN-27 prédit au sujet des triplets vérifiant $ q \geq 1.4 $.
\end{problem}

\noindent Les triplets vérifiant $ c > \operatorname{rad}(abc) $ ne sont pas
rares --- il y en a une infinité, et $ 1 + 8 = 9 $ est le plus petit --- de sorte que la
conjecture ne peut pas simplement les interdire ; tout le contenu tient dans l'exposant
$ 1 + \varepsilon $. Faire quatre exemples à la main rend cette distinction physique. Le
détenteur du record, trouvé par Éric Reyssat, est $ 2 + 3^{10} \cdot 109 = 23^{5} $, de
qualité environ $ 1.6299 $, et aucun triplet de qualité $ 2 $ ou plus n'a jamais été
trouvé --- la conjecture abc dit que, pour chaque $ \varepsilon > 0 $ fixé, seul un nombre
fini de triplets dépassent la qualité $ 1 + \varepsilon $.

\begin{refs}
\item Contexte : Conjecture abc --- Wikipédia (\url{https://fr.wikipedia.org/wiki/Conjecture_abc})
\item Contexte : Radical d'un entier --- Wikipédia (\url{https://fr.wikipedia.org/wiki/Radical_d%27un_entier})
\item Voir aussi : Théorie des nombres (\url{pure-number-theory.html})
\end{refs}

\bigskip

\begin{problem}[{Compter les points entiers dans un disque --- Licence, short}]
\textbf{OPEN-50.} \textbf{Problème.} Soit $ N(r) $ le nombre de points entiers
$ (x,y) \in \mathbb{Z}^2 $ vérifiant $ x^2 + y^2 \leq r^2 $. Calculez $ N(r) $ pour
$ r = 10, 20, 30, 50, 100 $, tabulez l'erreur $ E(r) = N(r) - \pi r^2 $, et décidez au
vu de vos données si $ E(r) $ semble plutôt de taille $ r^{1/2} $ ou de taille
$ r $.
\end{problem}

\noindent C'est le problème du cercle de Gauss, ouvert depuis que Gauss l'a
étudié vers 1800. Écrire $ N(r) = \pi r^2 + E(r) $ est trivial ; la question est de
savoir à quel point $ E(r) $ doit être petit. L'argument de Gauss lui-même donne
$ E(r) = O(r) $ --- chaque carré du bord est compté ou manqué --- et la conjecture est
$ E(r) = O(r^{1/2 + \varepsilon}) $, tandis que Hardy et Landau ont démontré en 1915 que
$ E(r) = O(r^{1/2}) $ est faux tel quel. La meilleure borne supérieure connue en 2026
est l'exposant de Huxley, $ 131/208 = 0.6298\ldots $ (2000) ; l'écart entre ce qui est
démontré et ce qui est cru reste donc ouvert. Cinq points de données ne trancheront pas,
mais ils vous montreront pourquoi l'exposant est difficile à voir.

\begin{refs}
\item Contexte : Problème du cercle de Gauss --- Wikipédia (\url{https://fr.wikipedia.org/wiki/Probl%C3%A8me_du_cercle_de_Gauss})
\item Lecture : G. H. Hardy \& E. M. Wright, \textit{An Introduction to the Theory of Numbers}, chap. 18
\end{refs}

\bigskip

\begin{problem}[{La conjecture de Beal --- Recherche}]
\textbf{OPEN-51.} \textbf{Problème.} Supposons que $ A, B, C, x, y, z $ soient des
entiers strictement positifs avec $ x, y, z > 2 $ et
\[ A^{x} + B^{y} = C^{z} . \]
\textbf{Démontrez que $ A $, $ B $ et $ C $ ont un facteur premier commun} --- ou
exhibez un contre-exemple. Statut : ouvert. Un prix d'un million de dollars américains,
placé en fiducie auprès de l'American Mathematical Society, est offert pour une
démonstration ou un contre-exemple publiés dans une revue à comité de lecture.
\end{problem}

\noindent D. Andrew Beal, banquier de Dallas et théoricien des nombres autodidacte,
est parvenu à cet énoncé au début des années 1990 en expérimentant sur des généralisations
du dernier théorème de Fermat ; il l'a annoncé avec un prix en 1997, et le fonds a été porté
par étapes à un million de dollars en 2013. La conjecture contient le dernier théorème de
Fermat comme cas particulier $ x = y = z = n > 2 $, puisqu'un facteur commun pourrait
être simplifié --- de sorte que toute démonstration devra englober le théorème de Wiles, ce
qui est la mesure honnête de sa difficulté. Ce que l'on sait provient de la même machinerie
de modularité : certains triplets d'exposants comme $ (3,3,n) $ ont été réglés en
attachant des courbes de Frey à des solutions hypothétiques, et de vastes calculs n'ont
trouvé aucun contre-exemple à petites bases et petits exposants. L'hypothèse du facteur
commun est essentielle et facile à tester : $ 3^{3} + 6^{3} = 3^{5} $, où tout est
divisible par $ 3 $.

\begin{refs}
\item Contexte : Équation de Fermat généralisée (conjecture de Beal) --- Wikipédia (\url{https://fr.wikipedia.org/wiki/%C3%89quation_de_Fermat_g%C3%A9n%C3%A9ralis%C3%A9e})
\item Prix : AMS --- règlement et procédure du prix Beal (\url{https://www.ams.org/profession/prizes-awards/ams-supported/beal-prize-rules})
\item Article : R. Daniel Mauldin, \og{} A generalization of Fermat's last theorem: the Beal conjecture and prize problem \fg{}, \textit{Notices of the AMS} 44 (1997)
\item Voir aussi : Théorie des nombres (\url{pure-number-theory.html})
\end{refs}

\bigskip

\begin{problem}[{La conjecture de Fermat--Catalan --- Recherche}]
\textbf{OPEN-52.} \textbf{Problème.} Considérez l'équation
\[ a^{m} + b^{n} = c^{k} \]
en entiers strictement positifs, avec $ a, b, c $ deux à deux premiers entre eux et des
exposants vérifiant
\[ \frac{1}{m} + \frac{1}{n} + \frac{1}{k} < 1 . \]
\begin{enumerate}
\item[(a)] Vérifiez les solutions connues $ 1 + 2^{3} = 3^{2} $, $ 2^{5} + 7^{2} = 3^{4} $ et
$ 7^{3} + 13^{2} = 2^{9} $, et contrôlez dans chaque cas que la condition sur les
exposants est satisfaite.
\item[(b)] \textbf{Démontrez qu'il n'existe qu'un nombre fini de telles solutions.} Statut :
ouvert. Dix solutions sont connues, la plus grande faisant intervenir des nombres de sept
chiffres et plus, et aucune n'a été trouvée en plus de vingt ans de recherche.
\end{enumerate}

\end{problem}

\noindent La condition sur les exposants est tout l'enjeu :
$ 1/m + 1/n + 1/k $ supérieur, égal ou inférieur à $ 1 $ correspond aux géométries
sphérique, euclidienne et hyperbolique sur un orbifold associé, et ce n'est que dans le
cas hyperbolique que les solutions devraient être rares. Darmon et Granville ont démontré
en 1995, à l'aide du théorème de Faltings, que pour chaque triplet d'exposants \textit{fixé}
il n'y a qu'un nombre fini de solutions premières entre elles ; la conjecture demande la
finitude lorsque les exposants varient aussi, ce qui est bien plus fort et découlerait de
la conjecture abc d'OPEN-27. Des jalons voisins montrent combien le terrain est délicat :
la conjecture de Catalan --- selon laquelle $ 8 $ et $ 9 $ sont les seules puissances
parfaites consécutives --- a résisté de 1844 jusqu'à ce que Preda Mihăilescu la démontre en
2002, et le dernier théorème de Fermat est le sous-cas $ a^{n} + b^{n} = c^{n} $.

\begin{refs}
\item Contexte : Conjecture de Fermat-Catalan --- Wikipédia (\url{https://fr.wikipedia.org/wiki/Conjecture_de_Fermat-Catalan})
\item Contexte : Conjecture de Catalan --- Wikipédia (\url{https://fr.wikipedia.org/wiki/Conjecture_de_Catalan})
\item Article : H. Darmon \& A. Granville, \og{} On the equations $ z^m = F(x,y) $ and $ Ax^p + By^q = Cz^r $ \fg{}, \textit{Bulletin of the LMS} 27 (1995)
\end{refs}

\bigskip

\begin{problem}[{La conjecture de Sendov --- \textit{résolue en août 2026} --- Recherche}]
\textbf{OPEN-53.} \textbf{Problème (résolu tout récemment ; la vérification est encore en
cours de stabilisation).} Soit $ p $ un polynôme complexe de degré $ n \geq 2 $
dont toutes les racines appartiennent au disque unité fermé. \textbf{Montrez que, pour toute
racine $ a $ de $ p $, le disque fermé $ \{ z : |z - a| \leq 1 \} $ contient au
moins une racine de la dérivée $ p' $.} Statut au moment où ces lignes sont écrites :
pendant soixante ans la question est restée ouverte dans le cas général --- démontrée pour
les degrés $ n \leq 8 $ (Brown et Xiang, 1999) et, par Terence Tao en 2020, pour tout
$ n $ assez grand, avec un seuil ineffectif qui ne réglait aucun degré précis au-dessus
de $ 8 $. En août 2026, Lech Mazur a annoncé une démonstration assistée par ordinateur
du cas général, produite avec une aide massive de l'IA et vérifiée dans l'assistant de
preuve Lean ; Tao a ensuite passé plusieurs jours à digérer, contextualiser et simplifier
l'argument, et il rapporte qu'il démontre un énoncé plus fort qui règle aussi la
conjecture de Phelps--Rodriguez. Tenez cela pour \textbf{résolu, mais pas encore expertisé par
une revue} : la vérification en Lean est un indice fort, et le processus ordinaire
d'évaluation par les pairs n'est pas encore allé à son terme.
\end{problem}

\noindent Blagovest Sendov a proposé cette conjecture en 1959 ; Walter Hayman l'a
imprimée dans son recueil de problèmes de 1967 en l'attribuant par erreur à Ljubomir Iliev,
et le nom est resté des années. Le fait de fond est le théorème de Gauss--Lucas : les racines
de $ p' $ se trouvent toujours dans l'enveloppe convexe des racines de $ p $, et ne
peuvent donc pas s'échapper du disque unité. La question de Sendov est le raffinement local
--- non pas \og{} quelque part dans le disque \fg{}, mais \og{} à distance au plus un de \textit{chaque}
racine \fg{} --- et ce léger renforcement a consommé plus de soixante ans et quelque chose comme
une centaine d'articles. L'article de Tao de 2020 mérite encore d'être étudié comme pièce
de méthode : il raisonne par l'absurde à partir d'une suite de contre-exemples minimaux,
extrait un objet limite et l'analyse par la probabilité et la théorie du potentiel, au prix
d'un degré-seuil incalculable --- exemple d'une netteté inhabituelle de théorème vrai
\og{} à terme \fg{} sans être vrai \og{} à partir d'ici, explicitement \fg{}. Le dénouement de 2026 est un
jalon d'une autre nature, et cette entrée est laissée ici, parmi les problèmes ouverts
classiques, comme trace délibérée de la vitesse à laquelle le sol peut bouger : une
conjecture qui avait tenu le domaine en échec pendant six décennies a été refermée par un
argument produit par machine qu'un expert humain a ensuite dû passer des jours à apprendre
à lire. Quoi que vous concluiez sur ce que cela signifie pour les mathématiques, remarquez
que la vérification en Lean et le processus d'expertise sont deux choses différentes, et
qu'une seule des deux a eu lieu.

\begin{refs}
\item Contexte : Conjecture d'Iliev-Sendov --- Wikipédia (\url{https://fr.wikipedia.org/wiki/Conjecture_d%27Iliev-Sendov})
\item Contexte : Théorème de Gauss-Lucas --- Wikipédia (\url{https://fr.wikipedia.org/wiki/Th%C3%A9or%C3%A8me_de_Gauss-Lucas})
\item Article : Terence Tao, \og{} Sendov's conjecture for sufficiently high degree polynomials \fg{}, \textit{Acta Mathematica} 229 (2022), arXiv:2012.04125 (\url{https://arxiv.org/abs/2012.04125})
\item Lecture : Terence Tao, \og{} A digestion of the proof of Sendov's conjecture \fg{} (12 août 2026), terrytao.wordpress.com (\url{https://terrytao.wordpress.com/2026/08/12/a-digestion-of-the-proof-of-sendovs-conjecture/}) --- l'exposé accessible de l'argument de 2026
\item Article : Lech Mazur, \og{} A computer-assisted proof of Sendov's conjecture \fg{} (5 août 2026), ProofAtlas (PDF) (\url{https://www.proofatlas.ai/papers/sendov-conjecture/SENDOV_CONJECTURE_PROOF_AUGUST_5_2026.pdf}) ; la formalisation en Lean (\url{https://www.proofatlas.ai/formalizations/sendov-conjecture/}) est déposée à côté
\item Lecture : Q. I. Rahman \& G. Schmeisser, \textit{Analytic Theory of Polynomials} (2002)
\item Voir aussi : Analyse complexe (\url{pure-complex-analysis.html}) (le problème CA-31 traite les petits cas)
\end{refs}

\bigskip

\begin{problem}[{La conjecture de Singmaster --- Recherche}]
\textbf{OPEN-54.} \textbf{Problème.} Pour un entier $ t \geq 2 $, soit $ N(t) $ le nombre de couples
$ (n,k) $ tels que $ \binom{n}{k} = t $ --- autrement dit, le nombre de fois où $ t $
apparaît dans le triangle de Pascal.
\begin{enumerate}
\item[(a)] Trouvez les huit occurrences de $ 3003 $, le détenteur actuel du record, et montrez
que tout $ t \geq 2 $ apparaît au moins deux fois.
\item[(b)] \textbf{Démontrez que $ N(t) $ est majoré par une constante absolue.} Statut :
ouvert. Les meilleures bornes connues sont $ N(t) = O(\log t / \log\log t) $ (Abbott,
Erdős et Hanson, 1974) et le raffinement de Kane en 2007
\[ N(t) = O\!\left( \frac{(\log t)(\log\log\log t)}{(\log\log t)^{3}} \right) . \]
\end{enumerate}

\end{problem}

\noindent David Singmaster a demandé en 1971 si les entrées du triangle de Pascal
peuvent se répéter arbitrairement souvent, et a supposé que la vraie borne pourrait être
aussi petite que $ 8 $ ou $ 10 $ ; personne n'a trouvé de nombre apparaissant plus de
huit fois, et $ 3003 $ est le seul nombre connu à apparaître huit fois. Une infinité de
nombres apparaissent six fois, via une identité construite à partir des nombres de
Fibonacci que Singmaster a trouvée lui-même --- la réponse n'est donc certainement pas \og{} au
plus quatre \fg{}. La conjecture est un problème diophantien déguisé : chaque répétition est
une solution de $ \binom{n}{k} = \binom{m}{j} $, un système d'équations polynomiales, et
même une borne $ O(1) $ pour le cas particulier $ k = 2 $, $ j = 3 $ relève de la
théorie des points entiers sur les courbes. En 2021, Matomäki, Radziwiłł, Shao, Tao et
Teräväinen ont démontré la conjecture à l'intérieur du triangle, là où $ k $ n'est trop
proche ni de $ 0 $ ni de $ n $ --- laissant les fines régions de bord, c'est-à-dire
exactement là où vivent les répétitions connues.

\begin{refs}
\item Contexte : Conjecture de Singmaster --- Wikipédia (\url{https://fr.wikipedia.org/wiki/Conjecture_de_Singmaster})
\item Article : K. Matomäki, M. Radziwiłł, X. Shao, T. Tao, J. Teräväinen, \og{} Singmaster's conjecture in the interior of Pascal's triangle \fg{}, \textit{Quarterly Journal of Mathematics} 73 (2022), arXiv:2106.03335 (\url{https://arxiv.org/abs/2106.03335})
\item Lecture : Richard K. Guy, \textit{Unsolved Problems in Number Theory}
\item Voir aussi : Combinatoire (\url{pure-combinatorics.html})
\end{refs}

\bigskip

\begin{problem}[{La conjecture de Zaremba --- Recherche}]
\textbf{OPEN-55.} \textbf{Problème.} Tout rationnel $ b/d $ sous forme irréductible admet un
développement en fraction continue fini
\[ \frac{b}{d} = \cfrac{1}{a_1 + \cfrac{1}{a_2 + \cfrac{1}{\ddots + \cfrac{1}{a_r}}}}
= [0; a_1, a_2, \dots, a_r] , \]
à quotients partiels $ a_i $ entiers strictement positifs.
\begin{enumerate}
\item[(a)] Échauffement : pour $ d = 2, 3, \dots, 20 $, trouvez pour chaque $ d $ un $ b $
premier avec $ d $ dont tous les quotients partiels sont au plus $ 5 $.
\item[(b)] \textbf{Démontrez la conjecture de Zaremba : il existe une constante absolue $ A $ telle
que tout entier $ d \geq 2 $ admette un $ b $ avec $ \gcd(b,d) = 1 $ et tous les
quotients partiels de $ b/d $ au plus égaux à $ A $.} Statut : ouvert ;
$ A = 5 $ est la valeur conjecturée et elle s'accorde à tous les indices numériques.
\end{enumerate}

\end{problem}

\noindent Stanisław Zaremba a posé la question au début des années 1970 pour une
raison tout à fait pratique : dans la méthode du \og{} bon point de réseau \fg{} en intégration
numérique, la qualité d'une règle de réseau de module $ d $ est gouvernée par les
quotients partiels de $ b/d $, et il faut donc disposer de modules dont les fractions
continues sont sages. L'obstruction, c'est que les fractions continues à quotients partiels
bornés forment un ensemble de Cantor de dimension de Hausdorff fractionnaire --- les
dénominateurs qu'elles produisent sont clairsemés --- et exiger qu'elles atteignent
\textit{tout} entier, c'est demander à un ensemble mince d'être arithmétiquement complet.
Bourgain et Kontorovich ont réalisé la percée en 2014, en démontrant la conjecture pour un
ensemble de dénominateurs de densité un, grâce au développement d'une méthode du cercle
pour les semi-groupes minces de $ \mathrm{SL}_2(\mathbb{Z}) $ --- machinerie qui anime aussi
les travaux modernes sur les empilements de cercles apolloniens. Shinnyih Huang a ensuite
ramené la constante à la valeur conjecturée $ 5 $ pour un ensemble de densité un. Densité
un n'est pas la totalité : l'ensemble exceptionnel reste infini, et aucun $ d $ explicite
n'est connu pour mettre la conjecture en défaut.

\begin{refs}
\item Contexte : Quotients partiels bornés (conjecture de Zaremba) --- Wikipédia (en anglais) (\url{https://en.wikipedia.org/wiki/Restricted_partial_quotients})
\item Article : J. Bourgain \& A. Kontorovich, \og{} On Zaremba's conjecture \fg{}, \textit{Annals of Mathematics} 180 (2014), arXiv:1107.3776 (\url{https://arxiv.org/abs/1107.3776})
\item Article : Shinnyih Huang, \og{} An improvement to Zaremba's conjecture \fg{} (2013), arXiv:1310.3772 (\url{https://arxiv.org/abs/1310.3772})
\item Lecture : A. Ya. Khintchine, \textit{Continued Fractions}
\end{refs}

\bigskip

\begin{problem}[{Les irrationnels algébriques sont-ils normaux ? --- Recherche}]
\textbf{OPEN-56.} \textbf{Problème.} Un nombre réel est \textit{normal} en base $ b $ si,
pour tout $ k \geq 1 $, chacun des $ b^{k} $ blocs de $ k $ chiffres apparaît dans
son développement en base $ b $ avec une fréquence limite $ b^{-k} $. \textbf{Démontrez que
tout nombre algébrique irrationnel --- $ \sqrt{2} $, disons, ou $ 2^{1/3} $ --- est normal
dans toute base.} Statut : ouvert, et ouvert à une profondeur embarrassante : aucun
nombre algébrique irrationnel n'a été démontré normal dans quelque base que ce soit, et
l'on ne sait même pas si le développement décimal de $ \sqrt{2} $ contient une infinité
d'occurrences du chiffre $ 7 $.
\end{problem}

\noindent Émile Borel a démontré en 1909 que presque tout nombre réel est normal
dans toute base --- les nombres non normaux forment un ensemble de mesure nulle --- et en 1950
il a conjecturé que les irrationnels algébriques figurent parmi les nombres normaux. Le
gouffre entre ces deux phrases est tout l'objet de cette entrée : une propriété peut valoir
pour presque tous les nombres et rester invérifiable pour un seul nombre que vous savez
nommer. Des constructions existent bien --- Champernowne a montré en 1933 que
$ 0{,}123456789101112\ldots $ est normal en base dix --- mais ce sont des nombres bâtis
pour être normaux, non des nombres apparus pour d'autres raisons. Les vrais résultats
partiels mesurent l'éloignement de la frontière : Bailey, Borwein, Crandall et Pomerance
ont démontré qu'un irrationnel algébrique de degré $ d $ compte au moins de l'ordre de
$ N^{1/d} $ chiffres 1 parmi ses $ N $ premiers chiffres binaires, et Adamczewski et
Bugeaud ont démontré que le nombre de blocs distincts de longueur $ k $ dans le
développement croît plus vite que linéairement en $ k $, de sorte qu'aucun irrationnel
algébrique n'a un développement aussi simple qu'une suite automatique. Les deux résultats
sont immensément plus faibles que \og{} chaque chiffre apparaît un dixième du temps \fg{}.

\begin{refs}
\item Contexte : Nombre normal --- Wikipédia (\url{https://fr.wikipedia.org/wiki/Nombre_normal})
\item Contexte : Constante de Champernowne --- Wikipédia (\url{https://fr.wikipedia.org/wiki/Constante_de_Champernowne})
\item Article : B. Adamczewski \& Y. Bugeaud, \og{} On the complexity of algebraic numbers I. Expansions in integer bases \fg{}, \textit{Annals of Mathematics} 165 (2007)
\item Article : D. H. Bailey, J. M. Borwein, R. E. Crandall, C. Pomerance, \og{} On the binary expansions of algebraic numbers \fg{}, \textit{Journal de Théorie des Nombres de Bordeaux} 16 (2004)
\item Voir aussi : Analyse réelle (\url{pure-real-analysis.html}), Théorie des nombres (\url{pure-number-theory.html})
\end{refs}

\bigskip

\begin{problem}[{La conjecture de Cramér : quelle taille un écart entre nombres premiers peut-il atteindre ? --- Recherche}]
\textbf{OPEN-61.} \textbf{Problème.} Notons $ g_n = p_{n+1} - p_n $ l'écart qui suit le n-ième nombre
premier. Cramér a conjecturé en 1936 que
\[ g_n \;=\; O\!\left( (\log p_n)^2 \right), \]
et, sous sa forme fine, que $ \limsup_{n \to \infty} g_n / (\log p_n)^2 = 1 $.
\begin{enumerate}
\item[(a)] Échauffez-vous avec les faits classiques : démontrez que les écarts entre nombres
premiers ne sont pas bornés, en exhibant n entiers composés consécutifs pour tout n, et
déduisez du théorème des nombres premiers que l'écart moyen au voisinage de x vaut environ
$ \log x $.
\item[(b)] Situez la frontière des deux côtés. Énoncez la meilleure borne supérieure
inconditionnelle --- Baker, Harman et Pintz (2001) donnent $ g_n < p_n^{0.525} $ pour n
grand --- et la meilleure borne connue sous l'hypothèse de Riemann, celle de Cramér :
$ g_n = O(\sqrt{p_n} \log p_n) $. Dites ensuite précisément ce qu'il faudrait améliorer
pour atteindre $ (\log p_n)^2 $, et notez que même en supposant OPEN-01 vous en restez
exponentiellement loin.
\item[(c)] Lisez la minoration de 2014 due à Ford, Green, Konyagin, Maynard et Tao : pour toute
constante c, une infinité de n vérifient
\[ g_n \;>\; c \, \log p_n \, \frac{\log\log p_n \, \log\log\log\log p_n}{(\log\log\log p_n)^2}. \]
Identifiez le point où la machinerie qui produit de \textit{grands} écarts se sépare de celle
d'OPEN-24, qui en produit de \textit{petits}, alors même que les deux sont issues du même
crible.
\item[(d)] \textbf{Démontrez ou réfutez la conjecture de Cramér.} Statut : ouvert, et même la
constante conjecturée est contestée --- voir le contexte ci-dessous.
\end{enumerate}

\end{problem}

\noindent Le modèle de Cramér --- colorier chaque entier n comme premier
indépendamment avec probabilité $ 1/\log n $ --- est le plus influent des modèles faux de
la théorie des nombres. Il prédit la bonne réponse à un très grand nombre de questions, et
une réponse démontrablement fausse à certaines : Helmut Maier a démontré en 1985 que le
nombre de nombres premiers dans les intervalles courts de longueur $ (\log x)^A $ fluctue
plus que le modèle ne le permet, raison pour laquelle le raffinement de Granville prédit une
limite supérieure d'au moins $ 2e^{-\gamma} \approx 1.1229 $ plutôt que le 1 de Cramér.
La minoration de la question (c) a réglé un problème qu'Erdős avait doté de 10 000 dollars,
son plus gros prix annoncé ; deux démonstrations sont parues sur arXiv à un jour
d'intervalle en août 2014, l'une signée de quatre auteurs et l'autre de Maynard seul, avant
d'être fondues en un article commun. Les outils étaient les mêmes méthodes de crible qui
venaient de produire les écarts bornés entre nombres premiers --- la machine identique
écartant les nombres premiers et les resserrant. Pour l'échelle : le plus grand écart
maximal jamais calculé vaut 1854, juste après le nombre premier 101412319996363309069,
tandis que $ (\log p)^2 $ y vaut à peu près 2200.

\begin{refs}
\item Contexte : Conjecture de Cramér --- Wikipédia (\url{https://fr.wikipedia.org/wiki/Conjecture_de_Cram%C3%A9r})
\item Contexte : Écart entre nombres premiers --- Wikipédia (\url{https://fr.wikipedia.org/wiki/%C3%89cart_entre_nombres_premiers})
\item Article : K. Ford, B. Green, S. Konyagin, J. Maynard \& T. Tao, \og{} Long gaps between primes \fg{}, \textit{Journal of the AMS} 31 (2018), 65--105
\item Voir aussi : Théorie des nombres (\url{pure-number-theory.html})
\end{refs}

\bigskip

\begin{problem}[{La liste de Hilbert, passée en revue --- Recherche}]
\textbf{OPEN-62.} \textbf{Problème.} Le 8 août 1900, au deuxième Congrès international des mathématiciens à
Paris, David Hilbert a donné une conférence sur les problèmes qui, selon lui, allaient
façonner le siècle à venir. Il en a lu dix à voix haute ; la version imprimée en compte 23.
Ce problème est une revue de compte, et c'est la revue qui est l'exercice.
\begin{enumerate}
\item[(a)] Construisez le tableau vous-même. Pour chacun des 23, écrivez une phrase énonçant le
problème, puis classez-le comme résolu (en nommant l'auteur de la solution et l'année),
partiellement résolu, non résolu, ou trop vague pour avoir une valeur de vérité --- et
défendez par écrit toute classification que vous n'avez pas pu simplement aller
chercher.
\item[(b)] Trois entrées figurent sur cette page à part entière. Le huitième problème de Hilbert
contient l'hypothèse de Riemann (OPEN-01), la conjecture de Goldbach (OPEN-23) et la
conjecture des nombres premiers jumeaux (OPEN-24), tous encore ouverts ; une partie du
dix-huitième est la conjecture de Kepler, tombée en 2014 (OPEN-38) ; le dixième a reçu une
réponse négative en 1970, avec un reste toujours ouvert (OPEN-63). Pour chacun, énoncez
exactement quelle part de la question de Hilbert a reçu réponse et quelle part n'en a pas
reçu.
\item[(c)] Vient maintenant la partie qui n'est pas de la comptabilité. Choisissez un problème
qui s'est révélé mal posé --- le sixième, qui demande une axiomatisation de la physique, et
le vingt-troisième, qui demande le développement ultérieur du calcul des variations, sont
les candidats habituels --- et écrivez deux pages sur ce que \og{} résoudre \fg{} aurait bien pu
vouloir dire. Que vous apprend le taux d'échec d'une liste sur la valeur des listes de
problèmes ?
\item[(d)] Enfin, comparez. Placez la liste de Hilbert à côté de la liste des prix du millénaire
de 2000, en haut de cette page, et à côté du \og{} vingt-quatrième problème \fg{} que Rüdiger
Thiele a trouvé dans les carnets de Hilbert en 2000 --- une demande de critère de simplicité
pour les démonstrations, que Hilbert avait rédigée et jamais publiée. Quelle liste
vieillira le mieux selon vous, et sur quels indices ?
\end{enumerate}

\end{problem}

\noindent Statut : plusieurs des 23 restent ouverts, au premier chef le huitième,
le douzième et le seizième, tandis que le quatrième, le vingt-troisième et le
vingt-quatrième resté inédit sont généralement jugés trop vagues pour être tranchés. La
liste a fonctionné d'une manière que personne n'aurait pu concevoir : des domaines entiers
sont nés d'une seule entrée --- la théorie de la démonstration du deuxième, la théorie des
ensembles diophantiens du dixième, la théorie moderne de la transcendance du septième --- et
les mathématiciens se sont disputé l'honneur d'en refermer un. Elle montre aussi les
limites de la prophétie. Les plus grandes surprises du vingtième siècle --- les théorèmes
d'incomplétude de Gödel, l'ordinateur, la théorie des catégories, la classification des
groupes finis simples --- sont presque entièrement absentes d'une liste dressée par le
mathématicien le plus complètement informé de son temps. Il y a aussi une belle ironie dans
le calendrier : la phrase la plus citée de Hilbert, \og{} wir müssen wissen, wir werden wissen \fg{},
ne vient pas de Paris en 1900 mais d'une allocution radiophonique à Königsberg en septembre
1930, prononcée à un jour près de la première annonce publique de l'incomplétude par Gödel,
lors d'un colloque tenu dans cette même ville.

\begin{refs}
\item Contexte : Problèmes de Hilbert --- Wikipédia (\url{https://fr.wikipedia.org/wiki/Probl%C3%A8mes_de_Hilbert})
\item Contexte : Vingt-quatrième problème de Hilbert --- Wikipédia (en anglais) (\url{https://en.wikipedia.org/wiki/Hilbert%27s_twenty-fourth_problem})
\item Lecture : Felix E. Browder (dir.), \textit{Mathematical Developments Arising from Hilbert Problems}, AMS (1976)
\item Lecture : Benjamin H. Yandell, \textit{The Honors Class: Hilbert's Problems and Their Solvers} (2002)
\end{refs}

\bigskip

\begin{problem}[{Le dixième problème de Hilbert sur les rationnels --- Recherche, short}]
\textbf{OPEN-63.} \textbf{Problème.} Iouri Matiiassevitch, achevant les travaux de Martin
Davis, Hilary Putnam et Julia Robinson, a démontré en 1970 qu'aucun algorithme ne décide si
une équation polynomiale arbitraire à coefficients entiers admet une solution en
\textit{entiers}. \textbf{Tranchez la même question en remplaçant les entiers par les rationnels :
existe-t-il un algorithme qui détermine si une équation polynomiale à coefficients
rationnels admet une solution rationnelle ?} Statut : ouvert.
\end{problem}

\noindent La voie classique vers une réponse négative sur un anneau R consiste à
écrire une définition diophantienne de $\mathbb{Z}$ à l'intérieur de R, puis à invoquer le théorème de
1970 --- et pour $\mathbb{Q}$, c'est exactement ce que personne ne sait faire. Les conjectures de Barry
Mazur sur la topologie des points rationnels vont en sens inverse : si les points réels de
l'ensemble des solutions rationnelles d'une variété n'ont toujours qu'un nombre fini de
composantes connexes, alors $\mathbb{Z}$ n'est pas diophantien sur $\mathbb{Q}$ et cette voie est morte, de sorte
qu'une démonstration d'indécidabilité devrait avoir une allure complètement différente. Le
cas voisin vient de bouger. En décembre 2024, Peter Koymans et Carlo Pagano ont déposé une
démonstration de l'indécidabilité du problème sur l'anneau des entiers de \textit{tout} corps
de nombres, et en janvier 2025 Levent Alpöge, Manjul Bhargava, Wei Ho et leurs coauteurs
ont déposé une démonstration indépendante du même résultat par d'autres méthodes ; ces
travaux sont récents et encore en cours de stabilisation, et aucun ne touche $\mathbb{Q}$ lui-même.

\begin{refs}
\item Contexte : Dixième problème de Hilbert --- Wikipédia (\url{https://fr.wikipedia.org/wiki/Dixi%C3%A8me_probl%C3%A8me_de_Hilbert})
\item Contexte : Ensemble diophantien --- Wikipédia (\url{https://fr.wikipedia.org/wiki/Diophantien})
\item Lecture : Yuri Matiyasevich, \textit{Hilbert's Tenth Problem} (MIT Press, 1993)
\item Voir aussi : FND-33 (\url{pure-foundations.html#FND-33}) --- la mise en perspective ensembliste et décidabilité sur la page Fondements.
\end{refs}

\bigskip

\begin{problem}[{$ \pi + e $ est-il irrationnel ? --- Recherche, short}]
\textbf{OPEN-64.} \textbf{Problème.} Démontrez d'abord, en n'utilisant que la
transcendance de $\pi$ et de e ainsi que le fait que les nombres algébriques forment un corps,
qu'au moins l'un des nombres $ \pi + e $ et $ \pi e $ est transcendant. \textbf{Tranchez
ensuite la vraie question : $ \pi + e $ est-il irrationnel ?} Statut : ouvert --- et
l'on ne sait même pas si $ \pi + e $, $ \pi e $, $ \pi/e $, $ \pi^{\pi} $,
$ e^{e} $ ou $ \pi^{e} $ est seulement irrationnel, encore moins transcendant.
\end{problem}

\noindent Hermite a démontré la transcendance de e en 1873 et Lindemann celle de
$\pi$ en 1882, mettant fin à deux mille ans de tentatives de quadrature du cercle ; un siècle
et demi plus tard, personne ne sait décider de l'irrationalité de leur somme. L'argument en
deux lignes de la première phrase --- considérez le trinôme unitaire dont les racines sont $\pi$
et e --- constitue à peu près tout ce que l'on sait, ce qui en fait la démonstration la moins
coûteuse de la minceur persistante de la théorie. La conjecture de Schanuel, formulée dans
les années 1960, réglerait cette question et presque toutes les questions du même genre
d'un seul coup, puisqu'elle implique que $\pi$ et e sont algébriquement indépendants ; qu'une
unique conjecture énorme porte tout le sujet est en soi le diagnostic. Gardez cette entrée
à l'esprit chaque fois qu'un énoncé vous semble trop évident pour mériter une
démonstration.

\begin{refs}
\item Contexte : Nombre transcendant --- Wikipédia (\url{https://fr.wikipedia.org/wiki/Nombre_transcendant})
\item Contexte : Conjecture de Schanuel --- Wikipédia (\url{https://fr.wikipedia.org/wiki/Conjecture_de_Schanuel})
\item Contexte : Théorème de Lindemann-Weierstrass --- Wikipédia (\url{https://fr.wikipedia.org/wiki/Th%C3%A9or%C3%A8me_de_Lindemann-Weierstrass})
\item Lecture : Alan Baker, \textit{Transcendental Number Theory} (Cambridge, 1975)
\item Voir aussi : CAL-32 (\url{applied-calculus.html#CAL-32}) --- la page Analyse, où il figure aux côtés des autres constantes que l'analyse ne sait pas cerner.
\end{refs}

\bigskip


\section*{Récemment tombés}

\begin{problem}[{Le dernier théorème de Fermat (1637--1995) --- Master}]
\textbf{OPEN-37.} \textbf{Problème (tombé).} Aucun triplet d'entiers strictement positifs ne vérifie
\[ a^n + b^n = c^n \qquad \text{pour } n > 2. \]
Démontré par Andrew Wiles en 1995, une étape clé étant achevée conjointement avec Richard
Taylor, en établissant la modularité des courbes elliptiques semi-stables. Comprenez
l'architecture.
\begin{enumerate}
\item[(a)] Énoncez la conjecture de modularité de Taniyama--Shimura--Weil.
\item[(b)] Expliquez la construction de Frey et le théorème de Ribet, qui montrent ensemble qu'un
contre-exemple à Fermat produirait une courbe elliptique non modulaire.
\item[(c)] Démontrez vous-même le cas $ n = 4 $ par la méthode de descente infinie de Fermat
lui-même.
\end{enumerate}

\end{problem}

\noindent Fermat a écrit son affirmation dans la marge d'un livre vers 1637 ; trois
siècles et demi d'efforts ont produit des branches entières de la théorie algébrique des
nombres --- les idéaux, la théorie du corps de classes, et bien d'autres choses --- comme
sous-produits de l'échec. Wiles a travaillé dans un quasi-secret pendant sept ans, a
annoncé la démonstration à Cambridge en juin 1993, puis a dû survivre à la découverte d'une
lacune sérieuse à l'automne, qu'il a refermée avec Taylor en septembre 1994. La structure
de la question (b) est ce qu'il faut admirer : le problème n'a pas été résolu de front, il
a été transformé en un énoncé portant sur un objet complètement différent.

\begin{refs}
\item Contexte : Dernier théorème de Fermat --- Wikipédia (\url{https://fr.wikipedia.org/wiki/Dernier_th%C3%A9or%C3%A8me_de_Fermat})
\item Contexte : Théorème de modularité --- Wikipédia (\url{https://fr.wikipedia.org/wiki/Th%C3%A9or%C3%A8me_de_modularit%C3%A9})
\item Lecture : Simon Singh, \textit{Fermat's Enigma} (1997) --- publié en français sous le titre \textit{Le Dernier Théorème de Fermat}
\item Voir aussi : Théorie des nombres (\url{pure-number-theory.html}), L'art de la démonstration (\url{proofs.html})
\end{refs}

\bigskip

\begin{problem}[{La conjecture de Kepler, et la démonstration par ordinateur (1611--2014) --- Master}]
\textbf{OPEN-38.} \textbf{Problème (tombé).} Aucun empilement de sphères congruentes dans
$ \mathbb{R}^3 $ n'a une densité supérieure à $ \pi/\sqrt{18} \approx 0.74048 $ --- la
densité de la pile de l'épicier. Thomas Hales a annoncé une démonstration en 1998 ; les
rapporteurs ont indiqué qu'ils étaient \og{} certains à 99 \% \fg{} mais ne pouvaient pas vérifier
intégralement les calculs, si bien que Hales a dirigé le projet Flyspeck, qui a produit en
2014 une démonstration formelle complète vérifiée par machine. Lisez l'histoire, puis
répondez par écrit :
\begin{enumerate}
\item[(a)] Que signifie, pour une démonstration, être \textit{vérifiée} ?
\item[(b)] Une démonstration formelle que vous ne pouvez pas lire vous-même vous donne-t-elle une
connaissance ?
\end{enumerate}

\end{problem}

\noindent Kepler a énoncé la conjecture en 1611 dans un opuscule sur les flocons
de neige. Hilbert l'a incluse dans son dix-huitième problème. La démonstration finale a
exigé de vérifier par ordinateur des milliers de problèmes d'optimisation non linéaire, ce
qui a forcé la communauté mathématique à affronter une question authentiquement nouvelle
sur ce qui constitue une démonstration --- la même question que soulevait le théorème des
quatre couleurs en 1976, et qui est aujourd'hui centrale dans l'essor des assistants de
preuve comme Lean. L'empilement de sphères a aussi produit l'une des grandes surprises
récentes : voir l'entrée suivante.

\begin{refs}
\item Contexte : Conjecture de Kepler --- Wikipédia (\url{https://fr.wikipedia.org/wiki/Conjecture_de_Kepler})
\item Contexte : Projet Flyspeck --- Wikipédia (en anglais) (\url{https://en.wikipedia.org/wiki/Flyspeck})
\item Lecture : George Szpiro, \textit{Kepler's Conjecture} (2003)
\end{refs}

\bigskip

\begin{problem}[{L'empilement de sphères en dimensions 8 et 24 (2016) --- Master}]
\textbf{OPEN-39.} \textbf{Problème (tombé).} Le réseau $ E_8 $ donne l'empilement de
sphères le plus dense en dimension $ 8 $, et le réseau de Leech en dimension $ 24 $.
Démontré par Maryna Viazovska en 2016 (dimension 8), puis quelques jours plus tard en
dimension 24 avec Cohn, Kumar, Miller et Radchenko. Comprenez la méthode : la borne de
programmation linéaire de Cohn--Elkies ramène le problème à la construction d'une fonction
auxiliaire \og{} magique \fg{} à zéros prescrits, et Viazovska l'a bâtie à partir de formes
modulaires. Notez ensuite ce qui subsiste : \textbf{la dimension 3 a exigé la démonstration par
ordinateur de Hales, et les dimensions autres que 1, 2, 3, 8, 24 sont ouvertes.}
\end{problem}

\noindent C'est l'un des résultats les plus stupéfiants du siècle. La borne de
Cohn--Elkies était connue depuis 2003 et, numériquement, elle était \textit{terriblement}
proche de la vérité exactement en dimensions 8 et 24 --- tout le monde voyait bien que la
fonction magique devait exister, et personne ne savait l'écrire. Viazovska l'a trouvée, dans
un article de vingt-trois pages signé d'elle seule, en utilisant les formes modulaires d'une
manière que nul n'avait anticipée ; elle a reçu la médaille Fields en 2022. C'est la
démonstration moderne la plus claire qu'un problème bloqué depuis une décennie peut tomber
parce qu'une seule personne a eu la bonne idée.

\begin{refs}
\item Contexte : Empilement de sphères --- Wikipédia (\url{https://fr.wikipedia.org/wiki/Empilement_compact})
\item Article : Maryna Viazovska, \og{} The sphere packing problem in dimension 8 \fg{} (2016), arXiv:1603.04246 (\url{https://arxiv.org/abs/1603.04246})
\item Contexte : Maryna Viazovska --- Wikipédia (\url{https://fr.wikipedia.org/wiki/Maryna_Viazovska})
\item Voir aussi : Géométrie (\url{pure-geometry.html})
\end{refs}

\bigskip

\begin{problem}[{La monotuile apériodique : le \og{} chapeau \fg{} (2023) --- Licence}]
\textbf{OPEN-40.} \textbf{Problème (tombé).} Pendant des décennies, on ignorait si une seule tuile pouvait
paver le plan de façon uniquement apériodique --- un \og{} einstein \fg{}, de l'allemand
\textit{ein Stein}, une pierre. En mars 2023, David Smith, Joseph Samuel Myers, Craig Kaplan
et Chaim Goodman-Strauss ont exhibé le \og{} chapeau \fg{}, un polycerf-volant à 13 côtés qui fait
exactement cela ; la même année, ils ont produit le \og{} spectre \fg{}, qui n'a besoin d'aucune
réflexion.
\begin{enumerate}
\item[(a)] Découpez des copies du chapeau et pavez avec, à la main.
\item[(b)] Lisez la stratégie de démonstration et expliquez comment un système de substitution
exclut la périodicité.
\end{enumerate}

\end{problem}

\noindent Les pavages apériodiques ont commencé avec l'ensemble de 20 426 tuiles
de Berger, en 1966, ont été réduits par Penrose en 1974 à deux seulement, puis ont stagné à
deux pendant un demi-siècle. La résolution est venue en partie d'un amateur : David Smith
est un technicien de l'imprimerie à la retraite qui a trouvé le chapeau en découpant des
formes dans du carton et en jouant avec, avant de l'apporter à des mathématiciens qui ont
fourni la démonstration. Tout l'épisode --- la forme d'un amateur, une vérification assistée
par ordinateur, un résultat annoncé d'abord sur un blog --- est un portrait de la façon dont
les mathématiques se font aujourd'hui.

\begin{refs}
\item Contexte : Problème einstein --- Wikipédia (\url{https://fr.wikipedia.org/wiki/Probl%C3%A8me_einstein})
\item Article : Smith, Myers, Kaplan \& Goodman-Strauss, \og{} An aperiodic monotile \fg{} (2023), arXiv:2303.10798 (\url{https://arxiv.org/abs/2303.10798})
\item Voir aussi : Géométrie (\url{pure-geometry.html})
\end{refs}

\bigskip

\begin{problem}[{Le problème du canapé (2024, en cours d'évaluation) --- Recherche}]
\textbf{OPEN-41.} \textbf{Problème (résolu tout récemment, vérification complète en
attente).} Quelle est la plus grande aire d'une forme rigide que l'on puisse manœuvrer
autour d'un coin à angle droit dans un couloir de largeur unité ? Gerver a construit en
1992 un canapé d'aire $ \approx 2.2195 $ et l'a conjecturé optimal. En novembre 2024,
Jineon Baek a déposé une démonstration de 119 pages établissant que le canapé de Gerver est
bien le maximum. Statut en 2026 : \textbf{la démonstration est en cours d'évaluation par les
pairs et n'est pas encore parue dans une revue, même si l'accueil a été positif et que les
spécialistes la jugent très probablement correcte.} Pendant ce temps, la variante
\og{} ambidextre \fg{}, où le couloir tourne dans les deux sens, reste ouverte.
\end{problem}

\noindent Leo Moser a posé le problème formellement en 1966. Son charme tient à ce
que c'est un vrai problème de recherche que l'on peut expliquer en déménageant des meubles,
et sa difficulté à ce que la forme optimale n'est donnée par aucune formule simple --- le
canapé de Gerver est assemblé à partir de dix-huit arcs analytiques distincts. Tout le monde
s'attendait à ce que la démonstration finale soit assistée par ordinateur ; celle de Baek ne
l'est pas. Cette entrée est délibérément laissée dans l'état \og{} en attente \fg{} pour montrer à
quoi ressemble le statut honnête d'un résultat tout frais --- l'acceptation par la communauté
mathématique est un processus, et cela prend du temps.

\begin{refs}
\item Contexte : Problème du sofa --- Wikipédia (\url{https://fr.wikipedia.org/wiki/Probl%C3%A8me_du_sofa})
\item Article : Jineon Baek, \og{} Optimality of Gerver's sofa \fg{} (2024), arXiv:2411.19826 (\url{https://arxiv.org/abs/2411.19826})
\item Lecture : Quanta Magazine --- The Largest Sofa You Can Move Around a Corner (en anglais) (\url{https://www.quantamagazine.org/the-largest-sofa-you-can-move-around-a-corner-20250214/})
\end{refs}

\bigskip

\begin{problem}[{La conjecture de Kakeya en dimension trois (2025) --- Recherche}]
\textbf{OPEN-42.} \textbf{Problème (résolu tout récemment).} Un ensemble de Kakeya
contient un segment de droite de longueur unité dans chaque direction. \textbf{Démontrez que
tout ensemble de Kakeya de $ \mathbb{R}^3 $ est de dimension de Hausdorff et de
Minkowski égale à $ 3 $.} En février 2025, Hong Wang et Joshua Zahl ont déposé une
démonstration, reçue comme un jalon ; Larry Guth en a écrit une introduction explicative.
La conjecture demeure \textbf{ouverte en dimension $ n \geq 4 $}, et les conjectures
apparentées de restriction et de Bochner--Riesz en analyse harmonique sont ouvertes.
\end{problem}

\noindent Besicovitch a montré en 1919 qu'un ensemble de Kakeya peut être de
mesure nulle, ce qui est déjà saisissant --- on peut faire pivoter une aiguille dans toutes
les directions à l'intérieur d'un ensemble sans aire. La conjecture affirme qu'un tel
ensemble doit néanmoins être dimensionnellement plein, et elle se trouve au centre d'un
réseau de problèmes d'analyse harmonique, d'EDP et de combinatoire additive, ce qui
explique que Quanta ait qualifié la démonstration d'événement \og{} d'une fois par siècle \fg{}.
Prenez cette entrée comme rappel que la frontière bouge : c'était encore, il y a quelques
années, un problème ouvert sur des listes exactement comme celle-ci.

\begin{refs}
\item Contexte : Ensemble de Besicovitch (Kakeya) --- Wikipédia (\url{https://fr.wikipedia.org/wiki/Ensemble_de_Besicovitch})
\item Article : Larry Guth, \og{} Introduction to the proof of the Kakeya conjecture \fg{} (2025), arXiv:2505.07695 (\url{https://arxiv.org/abs/2505.07695})
\item Lecture : Terence Tao --- The three-dimensional Kakeya conjecture, after Wang and Zahl (en anglais) (\url{https://terrytao.wordpress.com/2025/02/25/the-three-dimensional-kakeya-conjecture-after-wang-and-zahl/})
\item Voir aussi : Analyse réelle et théorie de la mesure (\url{pure-real-analysis.html})
\end{refs}

\bigskip

\begin{problem}[{La conjecture de Duffin--Schaeffer (1941--2019) --- Recherche}]
\textbf{OPEN-65.} \textbf{Problème (tombé).} Avec quelle finesse un nombre réel typique peut-il être
approché par des fractions ? Fixez une fonction $ f \colon \mathbb{N} \to [0, \infty) $
et demandez pour quelles f presque tout réel $\alpha$ admet une infinité de fractions $ p/q $
irréductibles vérifiant
\[ \left| \alpha - \frac{p}{q} \right| \;<\; \frac{f(q)}{q}. \]
Duffin et Schaeffer ont conjecturé en 1941 que cela se produit pour presque tout $\alpha$
exactement lorsque $ \sum_{q} \varphi(q) f(q)/q $ diverge, $\varphi$ étant l'indicatrice
d'Euler. Démontré par Dimitris Koukoulopoulos et James Maynard en 2019, publié dans les
\textit{Annals of Mathematics} en 2020.
\begin{enumerate}
\item[(a)] Démontrez vous-même la moitié facile : si $ \sum_q \varphi(q) f(q)/q $ converge,
alors presque tout $\alpha$ n'admet qu'un nombre fini de telles approximations. (Borel--Cantelli,
après avoir compté combien de p admissibles il y a pour chaque q.)
\item[(b)] Démontrez le théorème de Khintchine de 1924, ou reconstituez-le à partir des
références : le même critère, la condition de primalité entre p et q étant supprimée et f
supposée décroissante. Repérez ensuite l'étape exacte où la monotonie est utilisée.
\item[(c)] Construisez le contre-exemple de Duffin et Schaeffer, qui explique pourquoi la
conjecture a dû être formulée : une fonction f portée par un ensemble mince de
dénominateurs, avec $ \sum_q f(q) $ divergente, pour laquelle presque aucun $\alpha$ n'est
approchable. Expliquez ce que cela montre sur le critère de Khintchine une fois la
monotonie retirée.
\item[(d)] Lisez la démonstration de 2019 et décrivez son dispositif central --- un graphe pondéré
sur les dénominateurs dont les arêtes enregistrent combien de facteurs premiers deux
dénominateurs partagent --- et expliquez en un paragraphe pourquoi le second lemme de
Borel--Cantelli ne peut pas être appliqué tel quel, c'est-à-dire pourquoi les événements
associés à des q différents sont corrélés et en quoi le facteur d'indicatrice est la
correction due à la primalité mutuelle.
\end{enumerate}

\end{problem}

\noindent C'était le problème central de l'approximation diophantienne métrique et
il a tenu soixante-dix-huit ans. Le théorème de Khintchine de 1924 avait donné le motif ---
pour f monotone, une seule série divergente décide de tout, et la loi du zéro-un de
Gallagher, en 1961, garantit que l'ensemble des $\alpha$ approchables est de mesure exactement 0
ou exactement 1, et rien entre les deux. Ce que Duffin et Schaeffer ont vu, c'est que
supprimer la monotonie brise l'indépendance sur laquelle l'argument repose en secret : les
événements attachés à deux dénominateurs se recouvrent largement lorsque ces dénominateurs
partagent beaucoup de facteurs premiers, et aucune précaution portant sur une seule série
ne peut le voir. La démonstration de Koukoulopoulos et Maynard transforme la structure de
corrélation elle-même en objet combinatoire et y fait tourner un argument de compression
subtil, raison pour laquelle l'article se lit comme de la combinatoire additive plutôt que
comme de l'analyse. Maynard a reçu la médaille Fields en 2022, la citation mentionnant ce
travail à côté de ses résultats sur les écarts entre nombres premiers (OPEN-24,
OPEN-61).

\begin{refs}
\item Contexte : Conjecture de Duffin-Schaeffer --- Wikipédia (\url{https://fr.wikipedia.org/wiki/Conjecture_de_Duffin-Schaeffer})
\item Contexte : Approximation diophantienne --- Wikipédia (\url{https://fr.wikipedia.org/wiki/Approximation_diophantienne})
\item Article : D. Koukoulopoulos \& J. Maynard, \og{} On the Duffin--Schaeffer conjecture \fg{}, \textit{Annals of Mathematics} 192 (2020), arXiv:1907.04593 (\url{https://arxiv.org/abs/1907.04593})
\item Lecture : Glyn Harman, \textit{Metric Number Theory} (Oxford, 1998)
\end{refs}

\bigskip

\begin{problem}[{La conjecture de Mertens, réfutée (1897--1985) --- Master, short}]
\textbf{OPEN-66.} \textbf{Problème (tombé).} Soit $ M(x) = \sum_{n \le x} \mu(n) $ la
fonction de Mertens. À l'aide de la représentation intégrale
\[ \frac{1}{\zeta(s)} \;=\; s \int_{1}^{\infty} \frac{M(x)}{x^{s+1}} \, dx,
\qquad \operatorname{Re}(s) > 1, \]
démontrez que la conjecture de Mertens --- selon laquelle $ |M(x)| 1 $ --- aurait impliqué l'hypothèse de Riemann d'OPEN-01. Énoncez ensuite avec soin
ce que la réfutation de 1985 fait, et ne fait pas, à l'hypothèse de Riemann.
\end{problem}

\noindent Stieltjes a annoncé dans une lettre de 1885 à Hermite qu'il savait
démontrer que $ M(x)/\sqrt{x} $ est borné ; les tables de Mertens de 1897 appuyaient la
forme plus fine, avec la constante 1 ; les deux avaient tort. Andrew Odlyzko et Herman te
Riele ont réfuté la conjecture en 1985 en utilisant l'algorithme de réduction de base de
réseau LLL pour traquer une quasi-résonance parmi les deux mille premiers zéros de $\zeta$,
établissant que $ M(x)/\sqrt{x} $ dépasse 1 une infinité de fois et descend sous $-$1 une
infinité de fois --- sans exhiber le moindre contre-exemple. On n'en connaît toujours aucun :
la meilleure borne disponible situe seulement le premier en dessous d'environ
$ e^{1.96 \times 10^{19}} $, et les constantes limites ont depuis été poussées au-delà de
1,82 et de $-$1,83. Pendant ce temps, l'énoncé plus faible
$ M(x) = O(x^{1/2 + \varepsilon}) $ est, lui, réellement équivalent à l'hypothèse de
Riemann, et il est ouvert. Voici l'avertissement le plus net de la page : une conjecture
confirmée pour toutes les valeurs que quiconque savait calculer, crue pendant un siècle,
assez forte pour impliquer la conjecture la plus célèbre des mathématiques --- et
fausse.

\begin{refs}
\item Contexte : Conjecture de Mertens --- Wikipédia (\url{https://fr.wikipedia.org/wiki/Conjecture_de_Mertens})
\item Contexte : Fonction de Mertens --- Wikipédia (\url{https://fr.wikipedia.org/wiki/Fonction_de_Mertens})
\item Article : A. M. Odlyzko \& H. J. J. te Riele, \og{} Disproof of the Mertens conjecture \fg{}, \textit{Journal für die reine und angewandte Mathematik} 357 (1985), 138--160
\end{refs}

\bigskip


\section*{Pour aller plus loin}


\vfill
\noindent\emph{Hors de la Caverne} --- des probl\`emes, pas de r\'eponses.
\end{document}
